% !TeX program = pdflatex
\documentclass[a4paper]{article}
\usepackage[fontsize=9pt]{scrextend}

% =========================================================
% 2 Seiten auf 1 Blatt – A4 quer
% =========================================================
\usepackage{pgfpages}
\pgfpagesuselayout{2 on 1}[a4paper,landscape,border shrink=5mm]

\usepackage[utf8]{inputenc}
\usepackage[T1]{fontenc}
\usepackage[ngerman]{babel}
\usepackage{lmodern}

\usepackage{amsmath,amssymb,mathtools}
\usepackage{physics}
\usepackage{siunitx}
\usepackage{graphicx}
\usepackage{qrcode}
\usepackage{xcolor}
\usepackage{enumitem}
\usepackage{tcolorbox}
\tcbuselibrary{most}
\usepackage{titlesec}
\usepackage[
a4paper,
left=12mm,
right=12mm,
top=15mm,
bottom=15mm
]{geometry}
\setlength{\columnsep}{6mm}

% =========================================================
% Farbsystem (konsistent, nicht dekorativ)
%   accent: zentrale Definitionen/Gleichungen
%   info: Merksätze
%   ok: Vorgehen/Schritte
%   warn: Checks/Fehler
% =========================================================
\definecolor{accent}{HTML}{1E5AA8}
\definecolor{info}{HTML}{0B4F6C}
\definecolor{ok}{HTML}{1B7F3A}
\definecolor{warn}{HTML}{B45309}
\definecolor{danger}{HTML}{991B1B}
\definecolor{soft2}{RGB}{245,245,250}

% Section-Styles
\titleformat{\section}{\large\bfseries\color{accent}}{\thesection}{0.6em}{}
\titleformat{\subsection}{\normalsize\bfseries\color{info}}{\thesubsection}{0.6em}{}
\titleformat{\subsubsection}{\normalsize\bfseries\color{black}}{\thesubsubsection}{0.6em}{}

% =========================================================
% tcolorbox Basis-Setup
% =========================================================
\tcbset{
  enhanced,
  arc=2pt,
  boxrule=0.6pt,
  left=4pt,right=4pt,top=3pt,bottom=3pt,
  before skip=4pt, after skip=4pt,
  fonttitle=\bfseries,
}

% =========================================================
% Semantische Hervorhebungen (sparsam einsetzen)
% =========================================================
\newcommand{\imp}[1]{{\color{accent}\bfseries #1}}
\newcommand{\term}[1]{{\color{info}\bfseries #1}}
\newcommand{\warntext}[1]{{\color{danger}\bfseries #1}}

% Gleichungen: dezent hinterlegt
\newcommand{\eqhl}[1]{\colorbox{accent!6}{\ensuremath{\displaystyle #1}}}

% =========================================================
% Boxen (semantische Farben)
% =========================================================
\newtcolorbox{merke}{
  title=Merke,
  colback=info!4,
  colframe=info!55,
  coltitle=info!90!black
}
\newtcolorbox{vorgehen}{
  title=Vorgehen,
  colback=ok!4,
  colframe=ok!55,
  coltitle=ok!90!black
}
\newtcolorbox{schema}{
  title=Schema ,
  colback=accent!4,
  colframe=accent!55,
  coltitle=accent!90!black
}
\newtcolorbox{checkitem}{
  title=Prüfungs-Check,
  colback=warn!6,
  colframe=warn!60,
  coltitle=warn!90!black
}
\newtcolorbox{abwandlung}{
  title=Typische Abwandlungen,
  colback=warn!4,
  colframe=warn!55,
  coltitle=warn!90!black
}
\newtcolorbox{taktik}{
  title=Prüfungstaktik,
  colback=danger!4,
  colframe=danger!55,
  coltitle=danger!90!black
}

% =========================================================
% Makros (kompakt, prüfungsorientiert)
% =========================================================
\renewcommand{\dd}{\mathrm{d}}
\renewcommand{\grad}{\nabla}
\newcommand{\divg}{\nabla\cdot}
\renewcommand{\curl}{\nabla\times}

\newcommand{\E}{\vec E}
\newcommand{\D}{\vec D}
\newcommand{\Hf}{\vec H}
\newcommand{\B}{\vec B}
\newcommand{\J}{\vec J}
\newcommand{\A}{\vec A}

\newcommand{\eps}{\varepsilon}
\newcommand{\muo}{\mu_0}
\newcommand{\epso}{\varepsilon_0}

\setlist[itemize]{nosep,leftmargin=1.2em}
\setlist[enumerate]{nosep,leftmargin=1.6em}
\author{Gian Marco Näf}
\title{Angewandter Elektromagnetismus -- Zusammenfassung}
\date{\today}
\usepackage{hyperref}
% =========================================================
% Korrekte Seitenzahlen im TOC bei "2 on 1" (physische Blätter zählen)
% =========================================================
\newcounter{sheetpage}
\setcounter{sheetpage}{1}

% Nach jeder ausgelieferten *geraden* logischen Seite: neues Blatt
\AddToHook{shipout/after}{%
  \ifodd\value{page}\else\stepcounter{sheetpage}\fi
}

% Anzeige und TOC sollen Blattnummern verwenden
\renewcommand{\thepage}{\arabic{sheetpage}}

\hypersetup{colorlinks=true, linkcolor=accent, urlcolor=accent}





\begin{document}
\maketitle
\begin{center}
\vspace{1em}

\qrcode[height=2.5cm]{https://github.com/Giamo08/Zusammenfassungen}

\vspace{0.3em}

\href{https://github.com/Giamo08/Zusammenfassungen}{\texttt{github.com/Giamo08/Zusammenfassungen}}
\end{center}

\setcounter{tocdepth}{2}
\tableofcontents



\newpage


% =========================================================
\section{Methodik: Systematisches Lösen elektromagnetischer Aufgaben}
% =========================================================

\begin{merke}
In Prüfungen scheitert man selten an Algebra. Man scheitert an \imp{falschem Modell}.
Die Lösung ist ein standardisiertes Vorgehen:
\[
\text{Regime} \rightarrow \text{Unbekannte} \rightarrow \text{PDE} \rightarrow \text{Randbedingungen} \rightarrow \text{Lösung} \rightarrow \text{Check}.
\]
\end{merke}

% ---------------------------------------------------------
\subsection{Was ist eine ``elektromagnetische Aufgabe'' strukturell?}
% ---------------------------------------------------------

Eine typische Prüfungsaufgabe gibt dir:
\begin{itemize}
  \item eine \imp{Geometrie} (endlich oder unendlich, Symmetrie oder nicht),
  \item \imp{Materialdaten} (z.B. $\eps$, $\mu$, $\sigma$),
  \item \imp{Quellen} (Ladungen $\rho$, Ströme $\J$, vorgeschriebene Potentiale/Spannungen),
  \item eine \imp{Anregung} (DC oder harmonisch mit $f$ bzw. $\omega$),
  \item gesuchte Größen (z.B. $\varphi$, $\E$, $\J$, $\A$, $\B$, $\Hf$, Leistung, Impedanz).
\end{itemize}

Die Aufgabe ist immer ein \imp{Randwertproblem}:
\begin{itemize}
  \item Unbekannte Funktion(en) im Raum,
  \item führende Differentialgleichung,
  \item Randbedingungen an Oberflächen/Grenzen,
  \item ggf. Bedingungen im Unendlichen oder an Symmetrieachsen.
\end{itemize}

\begin{schema}
\imp{Merksatz:} Erst wenn du die \emph{PDE + Randbedingungen} sauber hast, darfst du rechnen.
Vorher ist jede Rechnung Glücksspiel.
\end{schema}

% ---------------------------------------------------------
\subsection{Der 30-Sekunden-Entscheidungsbaum}
% ---------------------------------------------------------

\begin{vorgehen}
\begin{enumerate}
  \item \imp{Zeitverhalten identifizieren:}
  \begin{itemize}
    \item statisch/DC: $\partial/\partial t = 0$
    \item harmonisch: $\propto e^{j\omega t}$
    \item Wellen/Propagation explizit: Wellengleichung / Ausbreitung
  \end{itemize}

  \item \imp{Welche Physik dominiert?} (oft schon aus Text klar)
  \begin{itemize}
    \item Ladungen/Spannung/Dielektrikum $\Rightarrow$ elektrisch dominiert
    \item Ströme/Induktion/Leiter $\Rightarrow$ magnetisch dominiert
    \item Grenzflächen + Reflexion/Transmission $\Rightarrow$ gekoppelt (Welle)
  \end{itemize}

  \item \imp{Primärvariable wählen:} $\varphi$ oder $\A$ oder direkt $\E,\Hf$ (Welle)
  \item \imp{Symmetrie bestimmen} und Koordinatensystem wählen
  \item \imp{Randbedingungen formulieren} (physikalisch, nicht geraten)
  \item \imp{Lösen} (DGL/PDE)
  \item \imp{Sekundärgrößen ableiten} (z.B. $\E=-\nabla\varphi$, $\B=\nabla\times\A$, $\J=\sigma\E$)
  \item \imp{Checks} (Dimensionen, Grenzfälle, Richtung, Energiefluss)
\end{enumerate}
\end{vorgehen}

% ---------------------------------------------------------
\subsection{Regime-Klassifikation}
% ---------------------------------------------------------

\subsubsection{Elektrostatisch}
Kriterien (praktisch):
\begin{itemize}
  \item keine Zeitabhängigkeit,
  \item keine induzierten Wirbelfelder relevant,
  \item es geht um Potential/Spannung/Feld aus Ladungen oder Randpotentialen,
  \item Leiter werden oft als Equipotentialflächen behandelt.
\end{itemize}

Primärvariable: \imp{Skalarpotential} $\varphi$.

Feldableitung:
\[
\E = -\nabla\varphi.
\]

Führende Gleichung (typisch):
\begin{itemize}
  \item Laplace: $\Delta \varphi = 0$ (ladungsfreier Bereich)
  \item Poisson: $\Delta \varphi = -\rho/\eps$ (mit Volumenladung)
\end{itemize}

\begin{checkitem}
Wenn im Gebiet \imp{keine} freie Ladung angegeben ist, starte mit $\Delta\varphi=0$.
Wenn Ladungsdichte gegeben ist, starte mit Poisson.
\end{checkitem}

\subsubsection{Stationäre Strömungsanalyse (DC-Leitung)}
Kriterien:
\begin{itemize}
  \item keine Zeitabhängigkeit,
  \item elektrische Leitfähigkeit $\sigma$ ist relevant,
  \item gesucht sind meist Stromdichte $\J$, Feld $\E$, Potentialabfall (Widerstand).
\end{itemize}

Physik-Kette:
\[
\E = -\nabla\varphi,\qquad \J=\sigma\E,\qquad \nabla\cdot \J = 0.
\]

Führende Gleichung:
\[
\nabla\cdot(\sigma\nabla\varphi)=0
\quad\Rightarrow\quad
\Delta\varphi=0\ \text{falls }\sigma=\text{konstant}.
\]

\begin{merke}
\imp{Wichtig:} Mathematisch wieder Laplace, aber physikalisch ist $\varphi$ jetzt ein \imp{Leitungs-}Potential (Spannungsabfall), nicht ``Ladungspotential''.
\end{merke}

\subsubsection{Magnetostatik}
Kriterien:
\begin{itemize}
  \item $\partial/\partial t=0$,
  \item stationäre Ströme $\J$ treiben das Feld,
  \item gesucht sind $\B$, $\Hf$ und oft Vektorpotential $\A$.
\end{itemize}

Grundrelation:
\[
\B=\nabla\times \A.
\]
Dann wird (unter üblicher Eichwahl) eine Poisson-Gleichung für $\A$ gelöst:
\[
\Delta \A = -\mu \J.
\]

\begin{checkitem}
Wenn Geometrie und Stromrichtung eine eindeutige Symmetrie haben, ist $\A$ oft nur in \imp{eine} Richtung orientiert (massive Vereinfachung).
\end{checkitem}

\subsubsection{Magneto-Quasistatik (MQS)}
Kriterien:
\begin{itemize}
  \item harmonisch: Frequenz $f$ oder $\omega$ gegeben,
  \item Leiter mit endlicher $\sigma$,
  \item Induktion ist relevant, aber Wellencharakter wird vernachlässigt,
  \item typisch: Skin-Effekt, Feld- und Stromverdrängung.
\end{itemize}

Primärvariable: häufig $\A$ im Frequenzbereich.

Führende Struktur:
\[
\Delta \A - j\omega\mu\sigma\,\A = -\mu \J_q
\]
(Details kommen in späteren Kapiteln; hier nur als Regime-Marker).

\subsubsection{Elektrodynamik (Wellen)}
Kriterien:
\begin{itemize}
  \item Ausbreitung, Reflexion, Transmission, Impedanz, Wellengleichung,
  \item gekoppelte $\E$- und $\Hf$-Felder,
  \item Materialdaten als $\eps$, $\mu$ (und ggf. $\sigma$).
\end{itemize}

Primärvariable: oft direkt $\E$ und $\Hf$ (Phasoren).

\begin{merke}
Wellenprobleme werden selten über $\varphi$ oder $\A$ gelöst, sondern über Feldansätze + Randbedingungen + Impedanzen.
\end{merke}

% ---------------------------------------------------------
\subsection{Primärvariable wählen: Regeln, die wirklich funktionieren}
% ---------------------------------------------------------

\begin{schema}
\imp{Regel 1:} Wenn $\curl \E = 0$ gilt, dann \imp{immer} $\E=-\nabla\varphi$ nutzen.\\
\imp{Regel 2:} Wenn $\divg \B=0$ und Ströme treiben, dann \imp{immer} $\B=\curl \A$ nutzen.\\
\imp{Regel 3:} Bei Wellen: Feldansatz (einlaufend/reflektiert/transmittiert) + Randbedingungen.
\end{schema}

Typische Prüfungsentscheidung:
\begin{itemize}
  \item Potentialaufgabe (Laplace/Poisson): wähle $\varphi$
  \item Strom-/Leitungsaufgabe: wähle $\varphi$, leite $\E,\J$ ab
  \item Magnetostatik/MQS: wähle $\A$, leite $\B,\Hf$ ab
  \item Wellen: wähle $\tilde{\E},\tilde{\Hf}$, nutze Impedanz
\end{itemize}

% ---------------------------------------------------------
\subsection{Symmetrieanalyse}
% ---------------------------------------------------------

\begin{merke}
Symmetrie ist kein ``Nice-to-have''. Symmetrie ist das Werkzeug, das dir Komponenten und Variablen eliminiert.
\end{merke}

Fragenkatalog:
\begin{itemize}
  \item Welche Transformation lässt das Problem unverändert? (Rotation, Spiegelung, Translation)
  \item Welche Koordinate(n) können dann \imp{nicht} im Ergebnis vorkommen?
  \item Welche Feldkomponenten würden die Symmetrie verletzen?
\end{itemize}

Konsequenzen:
\begin{itemize}
  \item $\varphi=\varphi(r)$ statt $\varphi(r,\theta,\phi)$
  \item $\A=\A_z(r)\,\vec e_z$ statt voller Vektor
  \item Reduktion PDE $\rightarrow$ ODE
\end{itemize}

\begin{checkitem}
Wenn du mehr als eine Raumvariable hast, prüfe nochmals die Symmetrie. In 80\% der Prüfungsfälle geht es auf 1 Variable runter.
\end{checkitem}

% ---------------------------------------------------------
\subsection{Randbedingungen: wie man sie korrekt herleitet}
% ---------------------------------------------------------

Randbedingungen sind die physikalische Kopplung zwischen Lösung und Geometrie.

\subsubsection{Typen von Randbedingungen}
\begin{itemize}
  \item \imp{Dirichlet}: Größe selbst vorgegeben (z.B. $\varphi = U$ auf Leiter)
  \item \imp{Neumann}: Ableitung/Fluss vorgegeben (z.B. normaler Fluss, Feldkomponente)
  \item \imp{Gemischt (Robin)}: Kombination, häufig in MQS/Wellen
  \item \imp{Stetigkeitsbedingungen}: an Materialgrenzen
\end{itemize}

\subsubsection{Physikalische Quellen für Randbedingungen}
\begin{itemize}
  \item \imp{Idealer Leiter (Elektrostatik):} $\E_t=0$ auf der Oberfläche, Leiter ist Equipotential $\Rightarrow \varphi=\text{konstant}$.
  \item \imp{Isolator-Oberfläche:} kein Strom durch die Oberfläche $\Rightarrow \J_n=0$.
  \item \imp{Materialgrenze:} bestimmte Komponenten sind stetig (Details in späteren Kapiteln).
  \item \imp{Unendlichkeit:} Felder/Potentiale müssen ``vernünftig'' abklingen, z.B. $\varphi\to 0$ für $r\to\infty$ wenn Referenz so gesetzt.
\end{itemize}

\begin{merke}
Randbedingungen werden \imp{immer} aus der Physik formuliert.
Erst danach werden sie mathematisch als Bedingungen an $\varphi$ oder $\A$ übersetzt.
\end{merke}

% ---------------------------------------------------------
\subsection{Standard-Arbeitsablauf}
% ---------------------------------------------------------

\begin{vorgehen}
\begin{enumerate}
  \item \imp{Regime festlegen} (statisch, stationär, MQS, Welle)
  \item \imp{Unbekannte wählen} ($\varphi$ oder $\A$ oder $\tilde{\E}$)
  \item \imp{Gleichung aufschreiben} (Laplace/Poisson/Diffusion/Welle)
  \item \imp{Symmetrie anwenden} $\Rightarrow$ reduzierte Gleichung
  \item \imp{Randbedingungen} (inkl. Unendlichkeit/Symmetrieachsen)
  \item \imp{Lösen} (Konstanten über Randbedingungen)
  \item \imp{Ableiten} der gesuchten Größen (z.B. $\E,\J,\B,\Hf$)
  \item \imp{Plausibilitätscheck} (Dimensionen, Grenzwerte, Richtung)
\end{enumerate}
\end{vorgehen}

% ---------------------------------------------------------
\subsection{Plausibilitätschecks}
% ---------------------------------------------------------

\begin{schema}
Checks sind nicht optional. Sie retten Punkte.
\end{schema}

\begin{checkitem}
\begin{itemize}
  \item \imp{Dimensionen:} stimmt Einheit jeder Größe?
  \item \imp{Grenzfälle:} $\omega\to 0$, $\sigma\to 0$, $r\to\infty$, $x\to 0$.
  \item \imp{Monotonie:} Potentiale ``springen'' nicht ohne Randbedingung.
  \item \imp{Energiefluss:} bei Wellen muss $\vec S=\E\times\Hf$ in Ausbreitungsrichtung zeigen.
  \item \imp{Stetigkeit:} an Materialgrenzen keine willkürlichen Diskontinuitäten.
\end{itemize}
\end{checkitem}

% ---------------------------------------------------------
\subsection{Typische Fehlerbilder}
% ---------------------------------------------------------

\begin{abwandlung}
\begin{itemize}
  \item \imp{Falsches Regime}: z.B. Wellengleichung verwendet obwohl quasistatisch genügt.
  \item \imp{Zu früh rechnen}: bevor Symmetrie/Randbedingungen klar sind.
  \item \imp{Randbedingungen geraten}: statt aus Leiter/Isolator/Übergang abgeleitet.
  \item \imp{Feld direkt gesucht}: obwohl Potential die saubere Primärvariable ist.
  \item \imp{Kein Grenzfallcheck}: führt zu unphysikalischen Ergebnissen ohne dass man es merkt.
\end{itemize}
\end{abwandlung}

\begin{taktik}
Wenn du im Kreis läufst: schreibe \imp{nur} die drei Dinge neu hin:
Regime, führende Gleichung, Randbedingungen. Danach geht es wieder vorwärts.
\end{taktik}

% =========================================================
\section{Maxwell-Gleichungen und ihre physikalischen Reduktionen}
% =========================================================

\begin{merke}
Alle elektromagnetischen Aufgaben basieren auf den Maxwell-Gleichungen.
In Prüfungen wird \imp{nie} der volle Satz gleichzeitig benutzt, sondern immer eine
\imp{physikalisch begründete Reduktion}.
\end{merke}

% ---------------------------------------------------------
\subsection{Maxwell-Gleichungen in Differentialform}
% ---------------------------------------------------------

In makroskopischer Form lauten die Maxwell-Gleichungen:

\begin{align}
\nabla \cdot \D &= \rho
&&\text{(Gauss-Gesetz elektrisch)} \label{eq:maxwell1} \\
\nabla \cdot \B &= 0
&&\text{(keine magnetischen Monopole)} \label{eq:maxwell2} \\
\nabla \times \E &= -\frac{\partial \B}{\partial t}
&&\text{(Induktionsgesetz)} \label{eq:maxwell3} \\
\nabla \times \Hf &= \J + \frac{\partial \D}{\partial t}
&&\text{(Amp\`ere-Maxwell)} \label{eq:maxwell4}
\end{align}

Diese Gleichungen gelten unabhängig von Geometrie, Koordinatensystem und Regime.

\subsection{Materialgleichungen}

Um die Maxwell-Gleichungen zu schließen, werden Materialgleichungen benötigt.
In Prüfungen wird fast immer ein lineares, isotropes Material angenommen:

\begin{align}
\D &= \eps \E \\
\B &= \mu \Hf \\
\J &= \sigma \E
\end{align}

\begin{checkitem}
\begin{itemize}
  \item $\eps$: elektrische Permittivität
  \item $\mu$: magnetische Permeabilität
  \item $\sigma$: elektrische Leitfähigkeit
\end{itemize}
Falls nichts anderes angegeben ist, sind diese Größen \imp{konstant}.
\end{checkitem}

% ---------------------------------------------------------
\subsection{Kontinuitätsgleichung (Ladungserhaltung)}
% ---------------------------------------------------------

Aus der Divergenz von Gleichung \eqref{eq:maxwell4} folgt:

\[
\nabla \cdot \J + \frac{\partial \rho}{\partial t} = 0
\]

Dies ist die \imp{Kontinuitätsgleichung} und beschreibt die lokale Ladungserhaltung.

\begin{merke}
Diese Gleichung gilt \imp{immer}, unabhängig vom Regime.
\end{merke}

\begin{checkitem}
Stationärer Fall:
\[
\frac{\partial \rho}{\partial t} = 0
\quad \Rightarrow \quad
\nabla \cdot \J = 0
\]
Diese Bedingung ist zentral für stationäre Leitungsprobleme.
\end{checkitem}

% ---------------------------------------------------------
\subsection{Reduktionsprinzip: Von Maxwell zur führenden PDE}
% ---------------------------------------------------------

\begin{schema}
\imp{Prüfungslogik:}  
Nicht ``alle Maxwell-Gleichungen lösen'', sondern:
\[
\text{Maxwell} \;\Rightarrow\; \text{Regime-Annahmen} \;\Rightarrow\; \text{eine führende Gleichung}.
\]
\end{schema}

Im Folgenden werden die wichtigsten Reduktionen systematisch hergeleitet.

% ---------------------------------------------------------
\subsection{Elektrostatisches Regime}
% ---------------------------------------------------------

\subsubsection{Physikalische Annahmen}

\begin{itemize}
  \item keine Zeitabhängigkeit: $\partial/\partial t = 0$
  \item keine induzierten Wirbelfelder
  \item elektrische Felder werden durch Ladungen und Randpotentiale erzeugt
\end{itemize}

Damit gilt aus \eqref{eq:maxwell3}:
\[
\nabla \times \E = 0
\]

\subsubsection{Einführung des Skalarpotentials}

Aus $\nabla \times \E = 0$ folgt:
\[
\E = -\nabla \varphi
\]

\begin{merke}
Das Skalarpotential existiert \imp{nur}, wenn $\nabla \times \E = 0$ gilt.
\end{merke}

\subsubsection{Poisson- und Laplace-Gleichung}

Einsetzen in \eqref{eq:maxwell1}:
\[
\nabla \cdot (\eps \nabla \varphi) = -\rho
\]

Für konstantes $\eps$:
\[
\Delta \varphi = -\frac{\rho}{\eps}
\quad \text{(Poisson-Gleichung)}
\]

Ladungsfreier Bereich ($\rho = 0$):
\[
\Delta \varphi = 0
\quad \text{(Laplace-Gleichung)}
\]

\begin{checkitem}
\begin{itemize}
  \item Poisson: Volumenladungen vorhanden
  \item Laplace: keine Volumenladungen
\end{itemize}
\end{checkitem}

% ---------------------------------------------------------
\subsection{Stationäre Strömungsanalyse}
% ---------------------------------------------------------

\subsubsection{Physikalische Annahmen}

\begin{itemize}
  \item zeitunabhängig
  \item elektrische Leitfähigkeit $\sigma > 0$
  \item keine Ladungsakkumulation im Volumen
\end{itemize}

\subsubsection{Feld--Strom-Beziehung}

Aus der Elektrostatik:
\[
\E = -\nabla \varphi
\]

Ohmsches Gesetz:
\[
\J = \sigma \E = -\sigma \nabla \varphi
\]

Mit der Kontinuitätsgleichung:
\[
\nabla \cdot \J = 0
\]

\subsubsection{Führende Gleichung}

\[
\nabla \cdot (\sigma \nabla \varphi) = 0
\]

Für konstantes $\sigma$:
\[
\Delta \varphi = 0
\]

\begin{merke}
Mathematisch identisch zur Laplace-Gleichung der Elektrostatik,
physikalisch jedoch ein \imp{Leitungsproblem}.
\end{merke}

% ---------------------------------------------------------
\subsection{Magnetostatisches Regime}
% ---------------------------------------------------------

\subsubsection{Physikalische Annahmen}

\begin{itemize}
  \item zeitunabhängig
  \item stationäre Ströme erzeugen magnetische Felder
  \item keine induktiven Effekte
\end{itemize}

Damit gilt:
\[
\nabla \times \Hf = \J,
\qquad
\nabla \cdot \B = 0
\]

\subsubsection{Vektorpotential}

Aus $\nabla \cdot \B = 0$ folgt:
\[
\B = \nabla \times \A
\]

\begin{merke}
Das Vektorpotential ist nicht eindeutig (Eichfreiheit).
In Prüfungen wird implizit eine Standard-Eichung verwendet.
\end{merke}

\subsubsection{Poisson-Gleichung für das Vektorpotential}

Mit $\B = \mu \Hf$ und $\nabla \times \Hf = \J$ ergibt sich:
\[
\nabla \times (\nabla \times \A) = \mu \J
\]

Unter der Coulomb-Eichung $\nabla \cdot \A = 0$:
\[
\Delta \A = -\mu \J
\]

\begin{checkitem}
Die Lösung dieser Gleichung liefert direkt $\A$.
Magnetische Feldgrößen folgen über Ableitungen.
\end{checkitem}

% ---------------------------------------------------------
\subsection{Harmonische Zeitabhängigkeit und Phasoren}
% ---------------------------------------------------------

\subsubsection{Phasor-Darstellung}

Für harmonische Anregung:
\[
f(t) = \Re\{\tilde f e^{j\omega t}\}
\]

Ableitungsregel:
\[
\frac{\partial}{\partial t} \rightarrow j\omega
\]

\begin{merke}
Im Frequenzbereich werden Differentialgleichungen zu algebraischen Gleichungen in $j\omega$.
\end{merke}

\subsubsection{Maxwell-Gleichungen im Frequenzbereich}

Beispiel:
\[
\nabla \times \tilde{\E} = -j\omega \tilde{\B}
\]
\[
\nabla \times \tilde{\Hf} = \tilde{\J} + j\omega \tilde{\D}
\]

Diese Form ist Grundlage für MQS- und Wellenprobleme.

% ---------------------------------------------------------
\subsection{Magneto-quasistatisches Regime (MQS)}
% ---------------------------------------------------------

\subsubsection{Physikalische Annahmen}

\begin{itemize}
  \item harmonische Zeitabhängigkeit
  \item Leiter mit endlicher Leitfähigkeit
  \item Induktion relevant
  \item Wellencharakter vernachlässigbar (Geometrie $\ll \lambda$)
\end{itemize}

\subsubsection{Diffusionsgleichung}

Im Frequenzbereich führt die MQS-Näherung typischerweise auf:
\[
\Delta \A - j\omega \mu \sigma \A = -\mu \J_q
\]

\begin{merke}
Diese Gleichung beschreibt die Diffusion elektromagnetischer Felder im Leiter
(Skin-Effekt).
\end{merke}

% ---------------------------------------------------------
\subsection{Elektrodynamisches Regime (Wellen)}
% ---------------------------------------------------------

\subsubsection{Physikalische Annahmen}

\begin{itemize}
  \item gekoppelte elektrische und magnetische Felder
  \item Ausbreitung im Raum
  \item Reflexion und Transmission an Materialgrenzen
\end{itemize}

\subsubsection{Wellengleichung}

Aus Maxwell ergibt sich für homogene Medien:
\[
\Delta \E - \mu\eps \frac{\partial^2 \E}{\partial t^2} = 0
\]

Im Frequenzbereich:
\[
\Delta \tilde{\E} + \omega^2 \mu \eps \tilde{\E} = 0
\]

\begin{merke}
In Wellenproblemen werden Felder direkt angesetzt
(einlaufend, reflektiert, transmittiert).
\end{merke}

% ---------------------------------------------------------
\subsection{Zusammenfassung und Prüfungscheck}
% ---------------------------------------------------------

\begin{checkitem}
\begin{itemize}
  \item Regime eindeutig bestimmt?
  \item Zeitabhängigkeit korrekt behandelt?
  \item Führende Gleichung klar identifiziert?
  \item Physikalische Annahmen explizit genannt?
\end{itemize}
\end{checkitem}

\begin{taktik}
Wenn du unsicher bist, welches Regime gilt:
prüfe, ob Induktion ($\partial \B/\partial t$) oder Verschiebungsstrom ($\partial \D/\partial t$)
vernachlässigt werden dürfen.
\end{taktik}
% =========================================================
\section{Feld- und Potentialkonzepte}
% =========================================================

\begin{merke}
In Prüfungen werden Felder selten direkt gelöst.
Die zentrale Technik ist der Übergang:
\[
\text{Potential} \;\Rightarrow\; \text{Feld} \;\Rightarrow\; \text{Fluss- und Stromgrößen}.
\]
Wer Potentiale korrekt einsetzt, reduziert Komplexität und Fehlerquote massiv.
\end{merke}

% ---------------------------------------------------------
\subsection{Warum Potentiale?}
% ---------------------------------------------------------

Elektromagnetische Felder unterliegen Nebenbedingungen:
\begin{itemize}
  \item $\nabla \times \E$ ist oft null oder vorgegeben,
  \item $\nabla \cdot \B = 0$ gilt immer,
  \item Randbedingungen sind meist einfacher für Potentiale formulierbar.
\end{itemize}

\begin{schema}
\imp{Prüfungsregel:}
\begin{itemize}
  \item Wenn $\nabla \times \E = 0$ gilt $\Rightarrow$ Skalarpotential verwenden.
  \item Wenn $\nabla \cdot \B = 0$ gilt $\Rightarrow$ Vektorpotential verwenden.
\end{itemize}
\end{schema}

Potentiale sind mathematische Hilfsgrößen, aber physikalisch konsistent,
da alle messbaren Größen daraus eindeutig ableitbar sind.

% ---------------------------------------------------------
\subsection{Elektrisches Skalarpotential}
% ---------------------------------------------------------

\subsubsection{Definition}

Gilt im elektrostatischen und stationären Regime:
\[
\nabla \times \E = 0
\quad \Rightarrow \quad
\E = -\nabla \varphi
\]

\begin{merke}
Das Skalarpotential existiert nur, wenn keine zeitabhängige Induktion wirkt.
Bei $\partial \B/\partial t \neq 0$ ist $\varphi$ allein nicht ausreichend.
\end{merke}

\subsubsection{Physikalische Bedeutung}

\begin{itemize}
  \item $\varphi$ ist die potentielle Energie pro Ladung
  \item Potentialdifferenzen sind messbar, absolute Potentiale nicht
  \item Referenzpotential ist frei wählbar
\end{itemize}

\begin{schema}
Nur Potentialdifferenzen haben physikalische Bedeutung.
Ein konstantes Offset von $\varphi$ ändert keine Feldgröße.
\end{schema}

% ---------------------------------------------------------
\subsection{Elektrisches Feld aus dem Potential}
% ---------------------------------------------------------

\subsubsection{Ableitung}

Allgemein:
\[
\E = -\nabla \varphi
\]

In kartesischen Koordinaten:
\[
\E =
\begin{pmatrix}
-\partial \varphi/\partial x \\
-\partial \varphi/\partial y \\
-\partial \varphi/\partial z
\end{pmatrix}
\]

\subsubsection{Interpretation}

\begin{itemize}
  \item Feldlinien zeigen in Richtung fallenden Potentials
  \item $\E$ ist senkrecht zu Equipotentialflächen
  \item hohe Potentialgradienten $\Rightarrow$ hohe Feldstärken
\end{itemize}

\begin{checkitem}
Wenn $\varphi$ konstant auf einer Fläche ist, ist $\E$ dort normal zur Fläche.
\end{checkitem}

% ---------------------------------------------------------
\subsection{Elektrische Flussdichte}
% ---------------------------------------------------------

\subsubsection{Definition}

\[
\D = \eps \E
\]

\begin{itemize}
  \item $\D$ enthält die Wirkung freier und gebundener Ladungen
  \item in linearen Medien proportional zu $\E$
\end{itemize}

\subsubsection{Bedeutung für Randbedingungen}

Aus dem Gaußschen Gesetz:
\[
\nabla \cdot \D = \rho
\]

\begin{schema}
\imp{Randbedingung:}
Die normale Komponente von $\D$ springt um die Flächenladungsdichte.
\end{schema}

% ---------------------------------------------------------
\subsection{Magnetisches Vektorpotential}
% ---------------------------------------------------------

\subsubsection{Motivation}

Da immer gilt:
\[
\nabla \cdot \B = 0
\]
existiert ein Vektorpotential $\A$ mit:
\[
\B = \nabla \times \A
\]

\begin{merke}
Das Vektorpotential ist keine physikalisch messbare Größe,
aber extrem nützlich für Rechnungen.
\end{merke}

% ---------------------------------------------------------
\subsection{Eichfreiheit des Vektorpotentials}
% ---------------------------------------------------------

Das Vektorpotential ist nicht eindeutig.
Für jede skalare Funktion $\psi$ gilt:
\[
\A' = \A + \nabla \psi
\quad \Rightarrow \quad
\nabla \times \A' = \B
\]

\begin{schema}
Diese Freiheit nennt man \imp{Eichfreiheit}.
In Prüfungen wird stillschweigend eine Standard-Eichung benutzt.
\end{schema}

Häufige Wahl:
\[
\nabla \cdot \A = 0
\quad \text{(Coulomb-Eichung)}
\]

% ---------------------------------------------------------
\subsection{Magnetisches Feld aus dem Vektorpotential}
% ---------------------------------------------------------

\subsubsection{Ableitung}

\[
\B = \nabla \times \A
\qquad
\Hf = \frac{1}{\mu}\B
\]

\begin{itemize}
  \item $\B$ beschreibt Flussdichte
  \item $\Hf$ ist die feldtreibende Größe in Maxwell
\end{itemize}

\begin{checkitem}
In linearen Medien ist $\B$ proportional zu $\Hf$.
\end{checkitem}

% ---------------------------------------------------------
\subsection{Zusammenhang zwischen $\A$ und $\E$ (zeitabhängig)}
% ---------------------------------------------------------

Im zeitabhängigen Fall gilt allgemein:
\[
\E = -\nabla \varphi - \frac{\partial \A}{\partial t}
\]

\begin{merke}
\begin{itemize}
  \item Elektrostatik: $\partial \A/\partial t = 0$
  \item MQS/Wellen: beide Terme relevant
\end{itemize}
\end{merke}

Diese Gleichung erklärt:
\begin{itemize}
  \item induzierte Wirbelfelder
  \item warum $\varphi$ allein nicht ausreicht
\end{itemize}

% ---------------------------------------------------------
\subsection{Stromdichte aus Feldgrößen}
% ---------------------------------------------------------

\subsubsection{Ohmsches Gesetz (lokal)}

\[
\J = \sigma \E
\]

In Kombination mit Potential:
\[
\J = -\sigma \nabla \varphi
\]

\begin{schema}
Diese Beziehung koppelt elektrische Felder an stationäre Strömung.
\end{schema}

% ---------------------------------------------------------
\subsection{Energie- und Leistungsbetrachtung}
% ---------------------------------------------------------

\subsubsection{Leistungsdichte}

\[
p = \E \cdot \J
\]

\begin{itemize}
  \item beschreibt lokale Verlustleistung
  \item relevant bei Leitungs- und MQS-Problemen
\end{itemize}

\subsubsection{Poynting-Vektor}

\[
\vec S = \E \times \Hf
\]

\begin{merke}
$\vec S$ beschreibt Richtung und Dichte des elektromagnetischen Energieflusses.
In Wellenproblemen ist $\vec S$ zentral.
\end{merke}

% ---------------------------------------------------------
\subsection{Typische Prüfungsfehler bei Potentialen}
% ---------------------------------------------------------

\begin{abwandlung}
\begin{itemize}
  \item $\varphi$ benutzt, obwohl $\partial \B/\partial t \neq 0$ relevant ist
  \item falsche Ableitungsrichtung bei $\E=-\nabla\varphi$
  \item Verwechslung von $\B$ und $\Hf$
  \item Eichfreiheit ignoriert oder missverstanden
\end{itemize}
\end{abwandlung}

% ---------------------------------------------------------
\subsection{Prüfungs-Check}
% ---------------------------------------------------------

\begin{checkitem}
\begin{itemize}
  \item Existiert ein Skalarpotential?
  \item Existiert ein Vektorpotential?
  \item Wurden Felder korrekt aus Potentialen abgeleitet?
  \item Sind Richtung und Dimension der Felder plausibel?
\end{itemize}
\end{checkitem}

\begin{taktik}
Wenn Felder kompliziert aussehen:
gehe einen Schritt zurück und prüfe,
ob ein Potential existiert, das die Rechnung vereinfacht.
\end{taktik}
% =========================================================
\section{Randwertprobleme und Randbedingungen}
% =========================================================

\begin{merke}
In elektromagnetischen Prüfungen entscheidet nicht die Differentialgleichung,
sondern die \imp{Randbedingungen}.
Falsche oder fehlende Randbedingungen führen selbst bei korrekter Rechnung
zu falschen Ergebnissen.
\end{merke}

% ---------------------------------------------------------
\subsection{Was ist ein Randwertproblem?}
% ---------------------------------------------------------

Ein Randwertproblem besteht aus drei Elementen:
\begin{itemize}
  \item einer \imp{Unbekannten} (z.\,B. $\varphi$, $\A$, $\E$, $\Hf$),
  \item einer \imp{Differentialgleichung} im Volumen,
  \item \imp{Randbedingungen} auf den Begrenzungen des Gebietes.
\end{itemize}

Ohne Randbedingungen existieren unendlich viele mathematische Lösungen.
Erst Randbedingungen selektieren die physikalisch richtige Lösung.

\begin{schema}
\imp{Prüfungsrealität:}
Die PDE ist meist Standard (Laplace, Poisson, Diffusion, Welle).
Die eigentliche Arbeit ist das saubere Formulieren der Randbedingungen.
\end{schema}

% ---------------------------------------------------------
\subsection{Physikalische Herkunft von Randbedingungen}
% ---------------------------------------------------------

Randbedingungen stammen \imp{immer} aus physikalischen Eigenschaften von:
\begin{itemize}
  \item idealen Leitern,
  \item realen Leitern,
  \item Isolatoren,
  \item Materialgrenzen,
  \item Symmetrieflächen,
  \item dem Verhalten im Unendlichen.
\end{itemize}

Sie werden \imp{nicht geraten} und \imp{nicht aus der Mathematik abgelesen}.

% ---------------------------------------------------------
\subsection{Typen von Randbedingungen}
% ---------------------------------------------------------

\subsubsection{Dirichlet-Randbedingungen}

Der Wert der Unbekannten ist vorgegeben:
\[
\varphi = \varphi_0 \quad \text{auf dem Rand}.
\]

Typische Quellen:
\begin{itemize}
  \item ideale Leiter (Elektrostatik): konstantes Potential,
  \item vorgegebene Spannungen,
  \item Referenzpotential (z.\,B. $\varphi=0$).
\end{itemize}

\begin{merke}
Dirichlet-Bedingungen fixieren das Potential selbst,
nicht dessen Ableitung.
\end{merke}

% ---------------------------------------------------------
\subsubsection{Neumann-Randbedingungen}

Die Ableitung der Unbekannten ist vorgegeben:
\[
\frac{\partial \varphi}{\partial n} = g
\]

Physikalische Bedeutung:
\begin{itemize}
  \item vorgegebene Feldkomponente,
  \item vorgegebener Fluss,
  \item kein Fluss durch eine Oberfläche.
\end{itemize}

Beispiel:
\[
\J_n = 0 \Rightarrow \frac{\partial \varphi}{\partial n} = 0
\quad \text{(Isolatoroberfläche)}
\]

% ---------------------------------------------------------
\subsubsection{Gemischte (Robin-)Randbedingungen}

Kombination aus Wert und Ableitung:
\[
a\varphi + b\frac{\partial \varphi}{\partial n} = c
\]

Typisch bei:
\begin{itemize}
  \item MQS-Problemen,
  \item Wellenproblemen,
  \item Übergang zu dissipativen Medien.
\end{itemize}

% ---------------------------------------------------------
\subsection{Randbedingungen in der Elektrostatik}
% ---------------------------------------------------------

\subsubsection{Idealer Leiter}

Im elektrostatischen Gleichgewicht gilt:
\begin{itemize}
  \item Feld im Leiterinneren: $\E = 0$
  \item Leiter ist Equipotentialfläche
\end{itemize}

Daraus folgt:
\[
\varphi = \text{konstant auf der Leiteroberfläche}.
\]

\begin{schema}
Elektrostatik + idealer Leiter $\Rightarrow$ Dirichlet-Bedingung.
\end{schema}

% ---------------------------------------------------------
\subsubsection{Grenze Leiter -- Isolator}

Tangentialkomponente des elektrischen Feldes:
\[
\E_t = 0
\]

Normalkomponente:
\[
\D_{n,2} - \D_{n,1} = \rho_s
\]

\begin{checkitem}
Ohne explizit gegebene Flächenladungsdichte ist $\D_n$ stetig.
\end{checkitem}

% ---------------------------------------------------------
\subsection{Randbedingungen bei stationärer Strömung}
% ---------------------------------------------------------

In leitenden Medien gilt:
\[
\J = -\sigma \nabla \varphi
\]

\subsubsection{Isolierte Oberfläche}

Kein Stromfluss:
\[
\J_n = 0
\quad \Rightarrow \quad
\frac{\partial \varphi}{\partial n} = 0
\]

\subsubsection{Strom- oder Spannungsanschluss}

\begin{itemize}
  \item Spannungsanschluss $\Rightarrow \varphi$ vorgegeben
  \item Stromanschluss $\Rightarrow \J_n$ vorgegeben
\end{itemize}

\begin{merke}
In Prüfungen werden meist Spannungen vorgegeben.
Ströme folgen aus der Lösung.
\end{merke}

% ---------------------------------------------------------
\subsection{Randbedingungen in der Magnetostatik}
% ---------------------------------------------------------

Grundrelationen:
\[
\nabla \cdot \B = 0,
\qquad
\nabla \times \Hf = \J
\]

\subsubsection{Normalkomponente von $\B$}

\[
\B_{n,2} - \B_{n,1} = 0
\]

Die Normalkomponente der magnetischen Flussdichte ist immer stetig.

\subsubsection{Tangentialkomponente von $\Hf$}

\[
\Hf_{t,2} - \Hf_{t,1} = \vec K
\]

$\vec K$: Oberflächenstromdichte.

\begin{checkitem}
Ohne Oberflächenstrom ist $\Hf_t$ stetig.
\end{checkitem}

% ---------------------------------------------------------
\subsection{Randbedingungen für das Vektorpotential}
% ---------------------------------------------------------

Da:
\[
\B = \nabla \times \A
\]

folgen Randbedingungen für $\A$ indirekt aus denen von $\B$ und $\Hf$.

Typische Konsequenzen:
\begin{itemize}
  \item Stetigkeit bestimmter Komponenten von $\A$
  \item Stetigkeit bestimmter Ableitungen von $\A$
\end{itemize}

\begin{merke}
In Prüfungen werden diese Bedingungen meist direkt vorgegeben oder implizit verwendet.
\end{merke}

% ---------------------------------------------------------
\subsection{Randbedingungen im Unendlichen}
% ---------------------------------------------------------

Für offene Gebiete gilt:
\[
\text{Lösung muss für } r \to \infty \text{ endlich bleiben}.
\]

Typische Konsequenzen:
\begin{itemize}
  \item keine exponentiell wachsenden Terme,
  \item keine logarithmischen Divergenzen ohne physikalische Quelle.
\end{itemize}

\begin{schema}
Unphysikalisch wachsende Lösungen werden verworfen,
auch wenn sie mathematisch korrekt sind.
\end{schema}

% ---------------------------------------------------------
\subsection{Symmetrie als Randbedingung}
% ---------------------------------------------------------

Symmetrieflächen wirken wie Randbedingungen.

Beispiele:
\begin{itemize}
  \item Spiegelebene $\Rightarrow$ normale Ableitung verschwindet
  \item Rotationssymmetrie $\Rightarrow$ keine Winkelabhängigkeit
\end{itemize}

\begin{merke}
Symmetrie reduziert das Gebiet und die Anzahl der Randbedingungen.
\end{merke}

% ---------------------------------------------------------
\subsection{Prüfungs-Check für Randbedingungen}
% ---------------------------------------------------------

\begin{checkitem}
\begin{itemize}
  \item Jede Randfläche physikalisch interpretiert?
  \item Typ der Randbedingung korrekt (Dirichlet/Neumann)?
  \item Stetigkeitsbedingungen an Materialgrenzen beachtet?
  \item Verhalten im Unendlichen berücksichtigt?
\end{itemize}
\end{checkitem}

\begin{taktik}
Wenn eine Lösung nicht eindeutig ist:
Es fehlt fast immer eine Randbedingung.
Suche nach Symmetrie oder einem Referenzpotential.
\end{taktik}

% =========================================================
\section{Symmetrie, Koordinatensysteme und Differentialoperatoren}
% =========================================================

\begin{merke}
Die Wahl von Symmetrie und Koordinatensystem entscheidet über die Schwierigkeit der Aufgabe.
In Prüfungen lassen sich die meisten Probleme auf \imp{eine Raumvariable} reduzieren.
\end{merke}

% ---------------------------------------------------------
\subsection{Warum Symmetrie vor Rechnung?}
% ---------------------------------------------------------

Symmetrie liefert \imp{harte Aussagen} über:
\begin{itemize}
  \item welche Feldkomponenten existieren dürfen,
  \item von welchen Koordinaten die Lösung abhängt,
  \item welche Terme der Differentialgleichung verschwinden.
\end{itemize}

\begin{schema}
\imp{Regel:}
Erst Symmetrie analysieren, dann Koordinatensystem wählen,
erst danach die Differentialgleichung aufschreiben.
\end{schema}

% ---------------------------------------------------------
\subsection{Typen von Symmetrie}
% ---------------------------------------------------------

\subsubsection{Translationssymmetrie}

Eigenschaften:
\begin{itemize}
  \item Problem ist invariant entlang einer Achse.
  \item Feldgrößen hängen nicht von dieser Koordinate ab.
\end{itemize}

Konsequenz:
\[
\frac{\partial}{\partial z}(\cdot) = 0
\]

Typisch bei:
\begin{itemize}
  \item unendlich langen Strukturen,
  \item Leitern oder Platten mit konstanter Geometrie.
\end{itemize}

% ---------------------------------------------------------
\subsubsection{Rotationssymmetrie}

Eigenschaften:
\begin{itemize}
  \item Invarianz unter Rotation um eine Achse.
  \item Keine Winkelabhängigkeit.
\end{itemize}

Konsequenz:
\[
\frac{\partial}{\partial \varphi}(\cdot) = 0
\]

Typisch bei:
\begin{itemize}
  \item Zylindern,
  \item koaxialen Strukturen,
  \item rotationssymmetrischen Leitern.
\end{itemize}

% ---------------------------------------------------------
\subsubsection{Spiegelsymmetrie}

Eigenschaften:
\begin{itemize}
  \item Invarianz unter Spiegelung an einer Ebene.
\end{itemize}

Konsequenzen:
\begin{itemize}
  \item Normalkomponente eines Vektors verschwindet auf der Spiegelebene
  \item oder Tangentialkomponente verschwindet, je nach Feldtyp.
\end{itemize}

\begin{merke}
Spiegelsymmetrie kann als \imp{Randbedingung} verwendet werden.
\end{merke}

% ---------------------------------------------------------
\subsection{Auswahl des Koordinatensystems}
% ---------------------------------------------------------

Das Koordinatensystem wird so gewählt, dass:
\begin{itemize}
  \item Symmetrieflächen mit Koordinatenflächen zusammenfallen,
  \item die Unbekannte möglichst nur von einer Koordinate abhängt.
\end{itemize}

\begin{schema}
\imp{Faustregel:}
\begin{itemize}
  \item Ebenen, Rechtecke $\Rightarrow$ kartesisch
  \item Zylinder, Drähte $\Rightarrow$ zylindrisch
  \item Kugeln, Punktquellen $\Rightarrow$ sphärisch
\end{itemize}
\end{schema}

% ---------------------------------------------------------
\subsection{Differentialoperatoren – Überblick}
% ---------------------------------------------------------

In der Elektromagnetik treten drei Operatoren ständig auf:
\begin{itemize}
  \item Gradient
  \item Divergenz
  \item Rotation
\end{itemize}

Diese Operatoren sind koordinatenabhängig.

% ---------------------------------------------------------
\subsection{Kartesische Koordinaten}
% ---------------------------------------------------------

Koordinaten: $(x,y,z)$

\subsubsection{Gradient}

\[
\nabla \varphi =
\begin{pmatrix}
\partial \varphi / \partial x \\
\partial \varphi / \partial y \\
\partial \varphi / \partial z
\end{pmatrix}
\]

\subsubsection{Divergenz}

Für einen Vektor $\vec F = (F_x,F_y,F_z)$:
\[
\nabla \cdot \vec F =
\frac{\partial F_x}{\partial x}
+\frac{\partial F_y}{\partial y}
+\frac{\partial F_z}{\partial z}
\]

\subsubsection{Rotation}

\[
\nabla \times \vec F =
\begin{vmatrix}
\vec e_x & \vec e_y & \vec e_z \\
\partial/\partial x & \partial/\partial y & \partial/\partial z \\
F_x & F_y & F_z
\end{vmatrix}
\]

\subsubsection{Laplace-Operator}

\[
\Delta \varphi =
\frac{\partial^2 \varphi}{\partial x^2}
+\frac{\partial^2 \varphi}{\partial y^2}
+\frac{\partial^2 \varphi}{\partial z^2}
\]

% ---------------------------------------------------------
\subsection{Zylindrische Koordinaten}
% ---------------------------------------------------------

Koordinaten: $(r,\varphi,z)$

\subsubsection{Gradient}

\[
\nabla \varphi =
\frac{\partial \varphi}{\partial r}\vec e_r
+\frac{1}{r}\frac{\partial \varphi}{\partial \varphi}\vec e_\varphi
+\frac{\partial \varphi}{\partial z}\vec e_z
\]

\subsubsection{Divergenz}

Für $\vec F = (F_r,F_\varphi,F_z)$:
\[
\nabla \cdot \vec F =
\frac{1}{r}\frac{\partial}{\partial r}(rF_r)
+\frac{1}{r}\frac{\partial F_\varphi}{\partial \varphi}
+\frac{\partial F_z}{\partial z}
\]

\subsubsection{Rotation}

\[
\nabla \times \vec F =
\begin{pmatrix}
\frac{1}{r}\frac{\partial F_z}{\partial \varphi}-\frac{\partial F_\varphi}{\partial z} \\
\frac{\partial F_r}{\partial z}-\frac{\partial F_z}{\partial r} \\
\frac{1}{r}\frac{\partial}{\partial r}(rF_\varphi)-\frac{1}{r}\frac{\partial F_r}{\partial \varphi}
\end{pmatrix}
\]

\subsubsection{Laplace-Operator}

\[
\Delta \varphi =
\frac{1}{r}\frac{\partial}{\partial r}
\left(r\frac{\partial \varphi}{\partial r}\right)
+\frac{1}{r^2}\frac{\partial^2 \varphi}{\partial \varphi^2}
+\frac{\partial^2 \varphi}{\partial z^2}
\]

\begin{checkitem}
Bei Rotationssymmetrie gilt $\partial/\partial \varphi = 0$.
\end{checkitem}

% ---------------------------------------------------------
\subsection{Sphärische Koordinaten}
% ---------------------------------------------------------

Koordinaten: $(r,\theta,\varphi)$

\subsubsection{Gradient}

\[
\nabla \varphi =
\frac{\partial \varphi}{\partial r}\vec e_r
+\frac{1}{r}\frac{\partial \varphi}{\partial \theta}\vec e_\theta
+\frac{1}{r\sin\theta}\frac{\partial \varphi}{\partial \varphi}\vec e_\varphi
\]

\subsubsection{Laplace-Operator}

\[
\Delta \varphi =
\frac{1}{r^2}\frac{\partial}{\partial r}
\left(r^2\frac{\partial \varphi}{\partial r}\right)
+\frac{1}{r^2\sin\theta}\frac{\partial}{\partial \theta}
\left(\sin\theta\frac{\partial \varphi}{\partial \theta}\right)
+\frac{1}{r^2\sin^2\theta}
\frac{\partial^2 \varphi}{\partial \varphi^2}
\]

\begin{merke}
Bei Kugelsymmetrie hängt $\varphi$ nur von $r$ ab.
\end{merke}

% ---------------------------------------------------------
\subsection{Reduktion durch Symmetrie – typische Fälle}
% ---------------------------------------------------------

\begin{schema}
\begin{itemize}
  \item Ebenensymmetrie $\Rightarrow$ 1D in kartesisch
  \item Rotationssymmetrie $\Rightarrow$ 1D in zylindrisch
  \item Kugelsymmetrie $\Rightarrow$ 1D in sphärisch
\end{itemize}
\end{schema}

In diesen Fällen reduziert sich die PDE auf eine gewöhnliche Differentialgleichung.

% ---------------------------------------------------------
\subsection{Prüfungs-Check}
% ---------------------------------------------------------

\begin{checkitem}
\begin{itemize}
  \item Wurde die Symmetrie explizit begründet?
  \item Passt das Koordinatensystem zur Geometrie?
  \item Wurden alle unnötigen Ableitungen eliminiert?
  \item Ist die reduzierte Gleichung minimal?
\end{itemize}
\end{checkitem}

\begin{taktik}
Wenn eine Rechnung explodiert:
Wechsle das Koordinatensystem oder prüfe erneut die Symmetrie.
Fast immer liegt dort der Fehler.
\end{taktik}

% =========================================================
% =========================================================
\section{Grundlegende Differentialgleichungen der Elektromagnetik}
% =========================================================

\begin{merke}
In EM-Prüfungen sind es fast immer dieselben Differentialgleichungen.
Der Schwierigkeitsgrad liegt nicht in der Form der Gleichung,
sondern in der korrekten Reduktion, Lösung und Interpretation.
\end{merke}

% ---------------------------------------------------------
\subsection{Einordnung: PDE vs. ODE}
% ---------------------------------------------------------

Die elektromagnetischen Grundgleichungen sind partielle Differentialgleichungen (PDEs).
Durch Symmetrie (Kapitel 5) reduzieren sie sich in Prüfungen fast immer auf
gewöhnliche Differentialgleichungen (ODEs).

\begin{schema}
\imp{Prüfungsrealität:}
\[
\text{PDE} \xrightarrow{\text{Symmetrie}} \text{ODE in einer Variablen}
\]
\end{schema}

Typische Unbekannte:
\begin{itemize}
  \item Skalarpotential $\varphi$
  \item Komponenten des Vektorpotentials $\A$
  \item Feldkomponenten im Wellenfall
\end{itemize}

% ---------------------------------------------------------
\subsection{Laplace-Gleichung}
% ---------------------------------------------------------

\subsubsection{Definition}

\[
\Delta u = 0
\]

Diese Gleichung beschreibt ladungs- bzw. quellenfreie Gebiete.

\subsubsection{Physikalische Bedeutung}

\begin{itemize}
  \item Elektrostatik ohne Volumenladungen
  \item Stationäre Strömung ohne Quellen/Senken
  \item Potentiale im quellenfreien Raum
\end{itemize}

\begin{merke}
Die Laplace-Gleichung ist rein geometrisch.
Die Lösung wird vollständig durch Randbedingungen bestimmt.
\end{merke}

% ---------------------------------------------------------
\subsubsection{Typische 1D-Reduktionen}
% ---------------------------------------------------------

\paragraph{Kartesisch (1D):}
\[
\frac{\dd^2 u}{\dd x^2} = 0
\quad \Rightarrow \quad
u(x) = Ax + B
\]

\paragraph{Zylindrisch (nur $r$-Abhängigkeit):}
\[
\frac{1}{r}\frac{\dd}{\dd r}
\left(r\frac{\dd u}{\dd r}\right)=0
\quad \Rightarrow \quad
u(r)=A\ln r + B
\]

\paragraph{Sphärisch (nur $r$-Abhängigkeit):}
\[
\frac{1}{r^2}\frac{\dd}{\dd r}
\left(r^2\frac{\dd u}{\dd r}\right)=0
\quad \Rightarrow \quad
u(r)=\frac{A}{r}+B
\]

\begin{checkitem}
Diese drei Lösungsformen sind absolute Prüfungsstandards.
\end{checkitem}

% ---------------------------------------------------------
\subsection{Poisson-Gleichung}
% ---------------------------------------------------------

\subsubsection{Definition}

\[
\Delta u = -f(\vec r)
\]

Mit Quelle $f(\vec r)$, z.\,B.:
\[
\Delta \varphi = -\frac{\rho}{\eps}
\]

\subsubsection{Physikalische Bedeutung}

\begin{itemize}
  \item Elektrostatik mit Volumenladungsdichte
  \item Magnetostatik mit Stromdichte als Quelle
\end{itemize}

\begin{merke}
Poisson = Laplace + Quelle.
Die homogene Lösung ist immer Teil der Gesamtlösung.
\end{merke}

% ---------------------------------------------------------
\subsubsection{Lösungsstruktur}
% ---------------------------------------------------------

Allgemein:
\[
u = u_\text{hom} + u_\text{part}
\]

\begin{itemize}
  \item $u_\text{hom}$ löst $\Delta u = 0$
  \item $u_\text{part}$ berücksichtigt die Quelle
\end{itemize}

In Prüfungen wird $u_\text{part}$ oft durch direktes Integrieren gefunden.

% ---------------------------------------------------------
\subsection{Gewöhnliche DGL zweiter Ordnung}
% ---------------------------------------------------------

Viele EM-Probleme führen auf ODEs der Form:
\[
\frac{\dd^2 u}{\dd x^2} + a\frac{\dd u}{\dd x} + b u = g(x)
\]

\subsubsection{Homogener Fall}

\[
\frac{\dd^2 u}{\dd x^2} + b u = 0
\]

Ansatz:
\[
u = e^{\lambda x}
\]

Charakteristische Gleichung:
\[
\lambda^2 + b = 0
\]

Lösungen:
\begin{itemize}
  \item $b>0$: sinusförmig
  \item $b<0$: exponentiell
\end{itemize}

\begin{schema}
Exponentielle Lösungen $\Rightarrow$ Dämpfung oder Wachstum.\\
Sinusförmige Lösungen $\Rightarrow$ Wellencharakter.
\end{schema}

% ---------------------------------------------------------
\subsection{Diffusionsgleichung (MQS)}
% ---------------------------------------------------------

\subsubsection{Form}

Im Frequenzbereich:
\[
\frac{\dd^2 u}{\dd x^2} - j\omega\mu\sigma\,u = 0
\]

\subsubsection{Physikalische Bedeutung}

\begin{itemize}
  \item elektromagnetische Diffusion in Leitern
  \item Skin-Effekt
  \item Eindringtiefe begrenzt
\end{itemize}

Lösungsansatz:
\[
u = e^{\pm \gamma x},
\qquad
\gamma = \sqrt{j\omega\mu\sigma}
\]

\begin{merke}
Reeller und imaginärer Teil von $\gamma$ bestimmen
Dämpfung und Phasenverschiebung.
\end{merke}

% ---------------------------------------------------------
\subsection{Wellengleichung}
% ---------------------------------------------------------

\subsubsection{Zeitbereich}

\[
\Delta u - \mu\eps \frac{\partial^2 u}{\partial t^2} = 0
\]

\subsubsection{Frequenzbereich}

\[
\Delta \tilde u + \omega^2 \mu\eps \tilde u = 0
\]

\subsubsection{1D-Lösung}

\[
\tilde u(x)=A e^{-jkx}+B e^{+jkx},
\qquad
k=\omega\sqrt{\mu\eps}
\]

\begin{schema}
\begin{itemize}
  \item $e^{-jkx}$: vorlaufende Welle
  \item $e^{+jkx}$: rücklaufende Welle
\end{itemize}
\end{schema}

% ---------------------------------------------------------
\subsection{Randverhalten und physikalische Selektion}
% ---------------------------------------------------------

Mathematisch zulässige Lösungen werden verworfen, wenn sie:
\begin{itemize}
  \item im Unendlichen divergieren,
  \item Symmetrien verletzen,
  \item physikalische Randbedingungen missachten.
\end{itemize}

\begin{merke}
Nicht jede mathematische Lösung ist physikalisch erlaubt.
\end{merke}

% ---------------------------------------------------------
\subsection{Zusammenfassung der Standardlösungen}
% ---------------------------------------------------------

\begin{schema}
\begin{itemize}
  \item Laplace $\Rightarrow$ linear / $\ln r$ / $1/r$
  \item Poisson $\Rightarrow$ homogen + partikular
  \item MQS-Diffusion $\Rightarrow$ gedämpfte Exponentialfunktionen
  \item Welle $\Rightarrow$ komplexe Exponentialfunktionen
\end{itemize}
\end{schema}

% ---------------------------------------------------------
\subsection{Prüfungs-Check}
% ---------------------------------------------------------

\begin{checkitem}
\begin{itemize}
  \item Wurde die richtige Gleichung gewählt?
  \item Wurde Symmetrie korrekt angewendet?
  \item Wurden unphysikalische Lösungen verworfen?
  \item Sind Grenzfälle konsistent?
\end{itemize}
\end{checkitem}

\begin{taktik}
Wenn du die Lösung nicht einordnen kannst:
Bestimme Grenzfälle ($x\to 0$, $x\to\infty$, $\omega\to 0$).
Das entlarvt falsche Terme sofort.
\end{taktik}
% =========================================================
\section{Frequenzbereich, komplexe Darstellung und Phasoren}
% =========================================================

\begin{merke}
Sobald eine Aufgabe harmonisch zeitabhängig ist, wird sie im Frequenzbereich gelöst.
Das ist kein mathematischer Trick, sondern eine systematische Vereinfachung.
\end{merke}

% ---------------------------------------------------------
\subsection{Harmonische Zeitabhängigkeit}
% ---------------------------------------------------------

Viele elektromagnetische Systeme werden mit sinusförmigen Signalen betrieben.
Eine zeitabhängige Größe $f(t)$ wird als harmonisch angenommen:
\[
f(t)=\Re\{\tilde f e^{j\omega t}\}
\]

Dabei gilt:
\begin{itemize}
  \item $\tilde f$: komplexe Amplitude (Phasor)
  \item $\omega=2\pi f$: Kreisfrequenz
\end{itemize}

\begin{schema}
\imp{Prüfungsregel:}
Im Frequenzbereich rechnet man nur mit $\tilde f$.
Die Zeitabhängigkeit wird am Ende wieder ergänzt.
\end{schema}

% ---------------------------------------------------------
\subsection{Ableitungsregeln im Frequenzbereich}
% ---------------------------------------------------------

Zeitliche Ableitungen werden zu algebraischen Faktoren:
\[
\frac{\partial}{\partial t} \rightarrow j\omega
\]

Beispiele:
\[
\frac{\partial f}{\partial t}
\rightarrow j\omega \tilde f
\qquad
\frac{\partial^2 f}{\partial t^2}
\rightarrow -\omega^2 \tilde f
\]

\begin{merke}
Dadurch werden zeitabhängige Differentialgleichungen zu rein räumlichen Gleichungen.
\end{merke}

% ---------------------------------------------------------
\subsection{Maxwell-Gleichungen im Frequenzbereich}
% ---------------------------------------------------------

Im Phasorbereich lauten die Maxwell-Gleichungen:

\begin{align}
\nabla \cdot \tilde{\D} &= \tilde{\rho} \\
\nabla \cdot \tilde{\B} &= 0 \\
\nabla \times \tilde{\E} &= -j\omega \tilde{\B} \\
\nabla \times \tilde{\Hf} &= \tilde{\J} + j\omega \tilde{\D}
\end{align}

\begin{merke}
Der Verschiebungsstrom $j\omega\tilde{\D}$ ist der Schlüssel für Wellenphänomene.
\end{merke}

% ---------------------------------------------------------
\subsection{Komplexe Materialgrößen}
% ---------------------------------------------------------

Leitfähige Medien werden im Frequenzbereich durch komplexe Materialparameter beschrieben.

\subsubsection{Komplexe Permittivität}

Ohmsches Gesetz:
\[
\tilde{\J}=\sigma\tilde{\E}
\]

Einsetzen in Amp\`ere-Maxwell:
\[
\nabla \times \tilde{\Hf} = j\omega\eps\tilde{\E} + \sigma\tilde{\E}
\]

Definition:
\[
\eps_\text{eff} = \eps - \frac{j\sigma}{\omega}
\]

\begin{merke}
Leitung und Verschiebungsstrom lassen sich zu einer komplexen Permittivität zusammenfassen.
\end{merke}

% ---------------------------------------------------------
\subsection{Komplexe Wellenzahl und Ausbreitung}
% ---------------------------------------------------------

Für homogene Medien ergibt sich die Wellengleichung:
\[
\Delta \tilde{\E} + k^2 \tilde{\E} = 0
\]

mit:
\[
k = \omega\sqrt{\mu\eps_\text{eff}}
\]

\begin{itemize}
  \item Re$(k)$: Phasenänderung
  \item Im$(k)$: Dämpfung
\end{itemize}

\begin{schema}
Gedämpfte Wellen $\Rightarrow$ komplexe Wellenzahl.
\end{schema}

% ---------------------------------------------------------
\subsection{Phasenlage und komplexe Amplituden}
% ---------------------------------------------------------

Eine komplexe Größe:
\[
\tilde f = |\tilde f| e^{j\phi}
\]

entspricht:
\[
f(t)=|\tilde f|\cos(\omega t + \phi)
\]

\begin{merke}
\begin{itemize}
  \item Betrag: Amplitude
  \item Argument: Phasenverschiebung
\end{itemize}
\end{merke}

% ---------------------------------------------------------
\subsection{Impedanzkonzept}
% ---------------------------------------------------------

\subsubsection{Definition}

Für ebene Wellen:
\[
Z=\frac{\tilde E}{\tilde H}
\]

Im homogenen Medium:
\[
Z=\sqrt{\frac{\mu}{\eps_\text{eff}}}
\]

\subsubsection{Physikalische Bedeutung}

\begin{itemize}
  \item Kopplung zwischen elektrischem und magnetischem Feld
  \item bestimmt Reflexion und Transmission an Grenzflächen
\end{itemize}

\begin{checkitem}
Bei verlustfreien Medien ist $Z$ reell.
Bei verlustbehafteten Medien ist $Z$ komplex.
\end{checkitem}

% ---------------------------------------------------------
\subsection{Frequenzbereich vs. Zeitbereich}
% ---------------------------------------------------------

\begin{schema}
\begin{itemize}
  \item Zeitbereich: anschaulich, aber rechenintensiv
  \item Frequenzbereich: algebraisch einfach, Standard in Prüfungen
\end{itemize}
\end{schema}

\begin{merke}
In EM-Prüfungen wird der Frequenzbereich bevorzugt.
\end{merke}

% ---------------------------------------------------------
\subsection{Typische Prüfungsfehler}
% ---------------------------------------------------------

\begin{abwandlung}
\begin{itemize}
  \item Vergessen des Faktors $j\omega$
  \item Vermischen von Zeit- und Frequenzbereich
  \item Fehlinterpretation komplexer Größen
  \item Betrag und Phase nicht getrennt betrachtet
\end{itemize}
\end{abwandlung}

% ---------------------------------------------------------
\subsection{Prüfungs-Check}
% ---------------------------------------------------------

\begin{checkitem}
\begin{itemize}
  \item Harmonik klar angenommen?
  \item Alle zeitlichen Ableitungen korrekt ersetzt?
  \item Komplexe Materialparameter konsistent verwendet?
  \item Ergebnis sinnvoll interpretiert (Amplitude + Phase)?
\end{itemize}
\end{checkitem}

\begin{taktik}
Wenn Unsicherheit besteht:
prüfe den Grenzfall $\omega\to 0$ (Übergang zu DC).
\end{taktik}
% =========================================================
\section{Magneto-Quasistatik, Diffusion und Skin-Effekt}
% =========================================================

\begin{merke}
Das magneto-quasistatische (MQS) Regime ist der Übergangsbereich
zwischen reiner Magnetostatik und voller Elektrodynamik.
In Prüfungen ist es eines der häufigsten Fehlerfelder.
\end{merke}

% ---------------------------------------------------------
\subsection{Was bedeutet ``quasistatisch''?}
% ---------------------------------------------------------

Quasistatisch bedeutet:
\begin{itemize}
  \item Zeitabhängigkeit ist vorhanden,
  \item Wellencharakter ist vernachlässigbar,
  \item Felder reagieren lokal verzögert, aber ohne räumliche Ausbreitung als Welle.
\end{itemize}

Typisches Kriterium:
\[
\text{charakteristische Länge } L \ll \lambda
\]

mit der Wellenlänge:
\[
\lambda = \frac{2\pi}{\omega\sqrt{\mu\eps}}
\]

\begin{schema}
\imp{Merksatz:}
MQS gilt, wenn die Geometrie klein gegenüber der Wellenlänge ist,
aber zeitliche Änderungen nicht ignoriert werden können.
\end{schema}

% ---------------------------------------------------------
\subsection{Physikalische Annahmen des MQS-Regimes}
% ---------------------------------------------------------

\begin{itemize}
  \item harmonische Zeitabhängigkeit ($e^{j\omega t}$),
  \item elektrische Leitfähigkeit $\sigma > 0$,
  \item Induktion relevant ($\partial \B/\partial t \neq 0$),
  \item Verschiebungsstrom vernachlässigbar:
  \[
  j\omega\eps \ll \sigma
  \]
\end{itemize}

\begin{merke}
Das Vernachlässigen des Verschiebungsstroms ist das zentrale MQS-Kriterium.
\end{merke}

% ---------------------------------------------------------
\subsection{Grundgleichungen im MQS-Regime}
% ---------------------------------------------------------

Aus Maxwell im Frequenzbereich folgt:

\[
\nabla \times \tilde{\E} = -j\omega \tilde{\B}
\]
\[
\nabla \times \tilde{\Hf} = \tilde{\J}
\]

Materialgleichungen:
\[
\tilde{\B} = \mu \tilde{\Hf},
\qquad
\tilde{\J} = \sigma \tilde{\E}
\]

% ---------------------------------------------------------
\subsection{Vektorpotential-Formulierung}
% ---------------------------------------------------------

Wie in der Magnetostatik wird eingeführt:
\[
\tilde{\B} = \nabla \times \tilde{\A}
\]

Das elektrische Feld ergibt sich aus:
\[
\tilde{\E} = -j\omega \tilde{\A}
\]

\begin{merke}
Im MQS-Regime dominiert der induktive Term.
Das Skalarpotential ist meist vernachlässigbar.
\end{merke}

% ---------------------------------------------------------
\subsection{Herleitung der Diffusionsgleichung}
% ---------------------------------------------------------

Einsetzen in $\tilde{\J}=\sigma\tilde{\E}$:
\[
\tilde{\J} = -j\omega\sigma \tilde{\A}
\]

Amp\`ere-Gesetz:
\[
\nabla \times \tilde{\Hf} = \tilde{\J}
\quad \Rightarrow \quad
\nabla \times \left(\frac{1}{\mu}\nabla \times \tilde{\A}\right)
= -j\omega\sigma \tilde{\A}
\]

Unter Standard-Eichung ($\nabla\cdot\tilde{\A}=0$):
\[
\Delta \tilde{\A} - j\omega\mu\sigma \tilde{\A} = 0
\]

\begin{schema}
Diese Gleichung beschreibt elektromagnetische Diffusion im Leiter.
\end{schema}

% ---------------------------------------------------------
\subsection{Eindringtiefe (Skin-Tiefe)}
% ---------------------------------------------------------

Für eindimensionale Probleme ergibt sich eine Lösung der Form:
\[
\tilde{\A}(x) = A_0 e^{-\gamma x}
\]

mit:
\[
\gamma = \sqrt{j\omega\mu\sigma}
\]

Die Skin-Tiefe ist definiert als:
\[
\delta = \sqrt{\frac{2}{\omega\mu\sigma}}
\]

\begin{merke}
\begin{itemize}
  \item Große $\omega$ $\Rightarrow$ kleine $\delta$
  \item Große $\sigma$ $\Rightarrow$ kleine $\delta$
\end{itemize}
\end{merke}

Physikalische Bedeutung:
\begin{itemize}
  \item Strom und Feld konzentrieren sich nahe der Oberfläche,
  \item das Leiterinnere wird abgeschirmt.
\end{itemize}

% ---------------------------------------------------------
\subsection{Feld- und Stromverteilung im Leiter}
% ---------------------------------------------------------

Aus $\tilde{\E}=-j\omega\tilde{\A}$ folgt:
\[
\tilde{\E}(x) \propto e^{-x/\delta}
\]

Damit:
\[
\tilde{\J}(x) = \sigma \tilde{\E}(x) \propto e^{-x/\delta}
\]

\begin{schema}
Der Skin-Effekt ist eine direkte Folge der Diffusionsgleichung.
\end{schema}

% ---------------------------------------------------------
\subsection{Grenzfälle}
% ---------------------------------------------------------

\subsubsection{Niederfrequenz-Grenze}

\[
\omega \to 0 \Rightarrow \delta \to \infty
\]

Konsequenz:
\begin{itemize}
  \item gleichmäßige Stromverteilung,
  \item Übergang zur Magnetostatik.
\end{itemize}

\subsubsection{Hochfrequenz-Grenze}

\[
\omega \to \infty \Rightarrow \delta \to 0
\]

Konsequenz:
\begin{itemize}
  \item Ströme nur noch an der Oberfläche,
  \item effektiver Querschnitt sinkt.
\end{itemize}

% ---------------------------------------------------------
\subsection{Leistungsverluste im MQS-Regime}
% ---------------------------------------------------------

Lokale Verlustleistungsdichte:
\[
p = \Re\{\tilde{\E}\cdot\tilde{\J}^*\}
\]

\begin{itemize}
  \item steigt mit Frequenz,
  \item steigt mit Leitfähigkeit,
  \item konzentriert sich nahe der Oberfläche.
\end{itemize}

\begin{merke}
Der Skin-Effekt erhöht den effektiven Widerstand bei AC-Betrieb.
\end{merke}

% ---------------------------------------------------------
\subsection{Typische Prüfungsfehler}
% ---------------------------------------------------------

\begin{abwandlung}
\begin{itemize}
  \item MQS verwendet, obwohl Wellenregime nötig wäre
  \item Verschiebungsstrom zu Unrecht vernachlässigt
  \item falsche Interpretation der Skin-Tiefe
  \item exponentielle Lösungen nicht physikalisch selektiert
\end{itemize}
\end{abwandlung}

% ---------------------------------------------------------
\subsection{Prüfungs-Check}
% ---------------------------------------------------------

\begin{checkitem}
\begin{itemize}
  \item Ist $L\ll\lambda$ erfüllt?
  \item Wurde der Verschiebungsstrom korrekt vernachlässigt?
  \item Diffusionsgleichung korrekt aufgestellt?
  \item Grenzfälle ($\omega\to 0$, $\omega\to\infty$) konsistent?
\end{itemize}
\end{checkitem}

\begin{taktik}
Wenn Unsicherheit besteht, ob MQS gilt:
vergleiche $\sigma$ und $\omega\eps$.
Der größere Term dominiert die Physik.
\end{taktik}
% =========================================================
\section{Elektromagnetische Wellen, Reflexion und Transmission}
% =========================================================

\begin{merke}
Sobald elektrische und magnetische Felder gekoppelt sind und sich räumlich ausbreiten,
befindet man sich im \imp{Wellenregime}.
In Prüfungen ist die saubere Trennung von einlaufender, reflektierter und transmittierter Welle entscheidend.
\end{merke}

% ---------------------------------------------------------
\subsection{Physikalische Einordnung des Wellenregimes}
% ---------------------------------------------------------

Das elektromagnetische Wellenregime gilt, wenn:
\begin{itemize}
  \item zeitliche Ableitungen \emph{nicht} vernachlässigt werden dürfen,
  \item der Verschiebungsstrom $j\omega\D$ relevant ist,
  \item Feldänderungen sich mit endlicher Geschwindigkeit ausbreiten.
\end{itemize}

Charakteristische Ausbreitungsgeschwindigkeit:
\[
v = \frac{1}{\sqrt{\mu\eps}}
\]

\begin{schema}
\imp{Abgrenzung:}
\begin{itemize}
  \item MQS: $L \ll \lambda$
  \item Welle: $L \sim \lambda$ oder $L > \lambda$
\end{itemize}
\end{schema}

% ---------------------------------------------------------
\subsection{Wellengleichung aus Maxwell}
% ---------------------------------------------------------

Aus den Maxwell-Gleichungen im homogenen Medium folgt für das elektrische Feld:
\[
\Delta \E - \mu\eps \frac{\partial^2 \E}{\partial t^2} = 0
\]

Im Frequenzbereich:
\[
\Delta \tilde{\E} + \omega^2 \mu\eps \tilde{\E} = 0
\]

Eine analoge Gleichung gilt für $\tilde{\Hf}$.

\begin{merke}
Elektrische und magnetische Felder erfüllen dieselbe Wellengleichung
und sind über Maxwell fest miteinander gekoppelt.
\end{merke}

% ---------------------------------------------------------
\subsection{Ebene Wellen}
% ---------------------------------------------------------

\subsubsection{Definition}

Eine ebene Welle ist eine Lösung der Form:
\[
\tilde{\E}(\vec r) = \tilde{\E}_0 e^{-j\vec k\cdot \vec r}
\]

mit:
\[
|\vec k| = k = \omega\sqrt{\mu\eps}
\]

\begin{itemize}
  \item $\vec k$: Ausbreitungsvektor
  \item $\tilde{\E}_0$: konstante Amplitude
\end{itemize}

% ---------------------------------------------------------
\subsubsection{Orthogonalitätsbeziehungen}

Für ebene Wellen gilt immer:
\[
\vec k \perp \tilde{\E},
\qquad
\vec k \perp \tilde{\Hf},
\qquad
\tilde{\E} \perp \tilde{\Hf}
\]

\begin{schema}
$\tilde{\E}$, $\tilde{\Hf}$ und $\vec k$ bilden ein rechtshändiges Tripel.
\end{schema}

% ---------------------------------------------------------
\subsection{Zusammenhang zwischen $\tilde{\E}$ und $\tilde{\Hf}$}
% ---------------------------------------------------------

Aus Maxwell folgt:
\[
\tilde{\Hf} = \frac{1}{Z}\, \hat k \times \tilde{\E}
\]

mit der Wellenimpedanz:
\[
Z = \sqrt{\frac{\mu}{\eps}}
\]

\begin{merke}
In ebenen Wellen sind $\tilde{\E}$ und $\tilde{\Hf}$ \imp{nicht unabhängig}.
Eine Feldgröße bestimmt die andere vollständig.
\end{merke}

% ---------------------------------------------------------
\subsection{Energiefluss elektromagnetischer Wellen}
% ---------------------------------------------------------

Der zeitgemittelte Energiefluss wird durch den Poynting-Vektor beschrieben:
\[
\vec S = \frac{1}{2}\Re\{\tilde{\E} \times \tilde{\Hf}^*\}
\]

\begin{itemize}
  \item Richtung von $\vec S$: Ausbreitungsrichtung
  \item Betrag von $\vec S$: Leistungsflussdichte
\end{itemize}

\begin{checkitem}
Bei verlustfreien Medien ist $\vec S$ parallel zu $\vec k$.
\end{checkitem}

% ---------------------------------------------------------
\subsection{Reflexion und Transmission an Materialgrenzen}
% ---------------------------------------------------------

Betrachtet wird eine ebene Welle, die senkrecht auf eine Grenzfläche trifft.

Es existieren drei Wellen:
\begin{itemize}
  \item einlaufende Welle
  \item reflektierte Welle
  \item transmittierte Welle
\end{itemize}

\begin{schema}
In Prüfungen immer explizit drei Felder ansetzen.
\end{schema}

% ---------------------------------------------------------
\subsection{Randbedingungen an der Grenzfläche}
% ---------------------------------------------------------

Aus Maxwell folgen für die tangentialen Komponenten:
\[
\E_{t,1} = \E_{t,2}
\qquad
\Hf_{t,1} = \Hf_{t,2}
\]

Diese Bedingungen koppeln die Amplituden der drei Wellen.

\begin{merke}
Reflexion und Transmission sind reine Konsequenzen der Randbedingungen.
\end{merke}

% ---------------------------------------------------------
\subsection{Reflexions- und Transmissionskoeffizienten}
% ---------------------------------------------------------

Für senkrechten Einfall ergeben sich:

\subsubsection{Reflexionskoeffizient}

\[
\Gamma = \frac{Z_2 - Z_1}{Z_2 + Z_1}
\]

\subsubsection{Transmissionskoeffizient}

\[
\tau = \frac{2Z_2}{Z_2 + Z_1}
\]

\begin{itemize}
  \item $Z_1$: Impedanz des einlaufenden Mediums
  \item $Z_2$: Impedanz des zweiten Mediums
\end{itemize}

\begin{checkitem}
\begin{itemize}
  \item $Z_2 = Z_1 \Rightarrow \Gamma = 0$ (keine Reflexion)
  \item $Z_2 \to \infty \Rightarrow \Gamma \to 1$ (vollständige Reflexion)
\end{itemize}
\end{checkitem}

% ---------------------------------------------------------
\subsection{Leistungsbetrachtung}
% ---------------------------------------------------------

Reflektierte und transmittierte Leistungen:
\[
R = |\Gamma|^2,
\qquad
T = \frac{Z_1}{Z_2}|\tau|^2
\]

Energieerhaltung:
\[
R + T = 1
\quad \text{(verlustfreie Medien)}
\]

\begin{merke}
Diese Gleichung ist ein wichtiger Plausibilitätscheck in Prüfungen.
\end{merke}

% ---------------------------------------------------------
\subsection{Verlustbehaftete Medien}
% ---------------------------------------------------------

Für leitfähige Medien ist:
\[
Z = \sqrt{\frac{\mu}{\eps - j\sigma/\omega}}
\]

Konsequenzen:
\begin{itemize}
  \item komplexe Impedanz,
  \item Dämpfung der transmittierten Welle,
  \item Phasenverschiebung zwischen $\E$ und $\Hf$.
\end{itemize}

\begin{schema}
Verluste $\Rightarrow$ Amplitudenabnahme und Energieabsorption.
\end{schema}

% ---------------------------------------------------------
\subsection{Typische Prüfungsfehler}
% ---------------------------------------------------------

\begin{abwandlung}
\begin{itemize}
  \item nur eine Welle angesetzt (Reflexion vergessen),
  \item falsche Richtung der reflektierten Welle,
  \item falsche Impedanz verwendet,
  \item keine Leistungsbilanz überprüft.
\end{itemize}
\end{abwandlung}

% ---------------------------------------------------------
\subsection{Prüfungs-Check}
% ---------------------------------------------------------

\begin{checkitem}
\begin{itemize}
  \item Wellenregime korrekt identifiziert?
  \item Drei Wellen explizit angesetzt?
  \item Randbedingungen korrekt angewendet?
  \item Energieerhaltung überprüft?
\end{itemize}
\end{checkitem}

\begin{taktik}
Wenn du bei Wellen unsicher bist:
zeichne $\vec k$, $\E$ und $\Hf$.
Die Geometrie klärt oft sofort die Vorzeichen und Richtungen.
\end{taktik}
% =========================================================
\addtocontents{toc}{\protect\setcounter{tocdepth}{1}}
\section{Mathematisches Cheatsheet}
% =========================================================

\begin{merke}
Dieses Kapitel deckt die gesamte Mathematik ab,
die in Elektromagnetik-Übungen und Prüfungen faktisch benutzt wird:
Ableitungen, Integrale, Differentialgleichungen, komplexe Rechnung,
Vektorrechnung, Randwertprobleme.
\end{merke}

% ---------------------------------------------------------
\subsection{Elementare Ableitungen und Integrale}
% ---------------------------------------------------------

\subsubsection{Potenzen und Logarithmen}

\[
\frac{\dd}{\dd x} x^n = n x^{n-1}
\qquad
\int x^n \dd x = \frac{x^{n+1}}{n+1}
\]

\[
\frac{\dd}{\dd x}\ln x = \frac{1}{x}
\qquad
\int \frac{1}{x}\dd x = \ln x
\]

\begin{checkitem}
$\ln r$ tritt systematisch bei zylindrischer Symmetrie auf.
\end{checkitem}

% ---------------------------------------------------------
\subsubsection{Exponentialfunktionen}

\[
\frac{\dd}{\dd x} e^{ax} = a e^{ax}
\qquad
\int e^{ax}\dd x = \frac{1}{a}e^{ax}
\]

\[
\frac{\dd}{\dd x} e^{-(\alpha+j\beta)x}
=-(\alpha+j\beta)e^{-(\alpha+j\beta)x}
\]

% ---------------------------------------------------------
\subsection{Partielle Ableitungen}
% ---------------------------------------------------------

\[
\frac{\partial}{\partial x}(xy^2)=y^2
\qquad
\frac{\partial}{\partial y}(xy^2)=2xy
\]

\begin{schema}
Bei partiellen Ableitungen werden alle anderen Variablen als konstant behandelt.
\end{schema}

% ---------------------------------------------------------
% ---------------------------------------------------------
% 10.X Trigonometrie + Hyperbelfunktionen + komplexe Darstellung (prüfungsrelevant)
%   Einfügen in Kapitel 10 (Mathematisches Cheatsheet)
% ---------------------------------------------------------
\subsection{Trigonometrie, Hyperbelfunktionen und komplexe Darstellung}

\begin{merke}
In EM tauchen \textbf{sin/cos} und \textbf{sinh/cosh} systematisch als Lösungen von linearen DGLs auf:
\[
u'' + k^2 u = 0 \Rightarrow \sin,\cos
\qquad\qquad
u'' - k^2 u = 0 \Rightarrow \sinh,\cosh
\]
Im Frequenzbereich sind trigonometrische Terme oft nur \textbf{Real-/Imaginärteil} komplexer Exponentialansätze.
\end{merke}

\subsubsection{Definitionen über Exponentialfunktionen }
\begin{schema}
\[
e^{jx} = \cos x + j\sin x
\qquad\Rightarrow\qquad
\cos x = \frac{e^{jx}+e^{-jx}}{2},\;\;
\sin x = \frac{e^{jx}-e^{-jx}}{2j}
\]
\[
\sinh x = \frac{e^{x}-e^{-x}}{2},\;\;
\cosh x = \frac{e^{x}+e^{-x}}{2},\;\;
\tanh x = \frac{\sinh x}{\cosh x}
\]
\end{schema}

\begin{checkitem}
\begin{itemize}
  \item \imp{Kernintuition:} $\sin/\cos$ sind \textbf{oszillierend}, $\sinh/\cosh$ sind \textbf{wachsend/abklingend}.
  \item \imp{Für Randverhalten:} exponentiell wachsende Terme werden oft durch Physik/Randbedingungen eliminiert.
\end{itemize}
\end{checkitem}

\subsubsection{Ableitungen}
\begin{schema}
\[
\frac{d}{dx}\sin x = \cos x
\qquad
\frac{d}{dx}\cos x = -\sin x
\qquad
\frac{d}{dx}\tan x = \sec^2 x = \frac{1}{\cos^2 x}
\]
\[
\frac{d}{dx}\sinh x = \cosh x
\qquad
\frac{d}{dx}\cosh x = \sinh x
\qquad
\frac{d}{dx}\tanh x = \sech^2 x = \frac{1}{\cosh^2 x}
\]
\end{schema}

\subsubsection{Integrale }
\begin{schema}
\[
\int \sin x\,dx = -\cos x + C
\qquad
\int \cos x\,dx = \sin x + C
\qquad
\int \sec^2 x\,dx = \tan x + C
\]
\[
\int \sinh x\,dx = \cosh x + C
\qquad
\int \cosh x\,dx = \sinh x + C
\qquad
\int \sech^2 x\,dx = \tanh x + C
\]
\end{schema}

\begin{schema}
\imp{Sehr häufig (Substitution):}
\[
\int e^{ax}\sin(bx)\,dx
=
\frac{e^{ax}}{a^2+b^2}\big(a\sin(bx)-b\cos(bx)\big)+C
\]
\[
\int e^{ax}\cos(bx)\,dx
=
\frac{e^{ax}}{a^2+b^2}\big(a\cos(bx)+b\sin(bx)\big)+C
\]
\end{schema}

\subsubsection{Identitäten }
\begin{schema}
\imp{Trigonometrie:}
\[
\sin^2 x+\cos^2 x = 1
\qquad
1+\tan^2 x = \sec^2 x
\]
\[
\cos(a\pm b)=\cos a\cos b \mp \sin a\sin b
\qquad
\sin(a\pm b)=\sin a\cos b \pm \cos a\sin b
\]
\end{schema}

\begin{schema}
\imp{Hyperbel:}
\[
\cosh^2 x-\sinh^2 x = 1
\qquad
1-\tanh^2 x = \sech^2 x
\]
\[
\cosh(a\pm b)=\cosh a\cosh b \pm \sinh a\sinh b
\qquad
\sinh(a\pm b)=\sinh a\cosh b \pm \cosh a\sinh b
\]
\end{schema}

\subsubsection{Komplexe Argumente und Beziehungen zwischen trig und hyperbel}
\begin{schema}
\imp{Zentrale Brücke:}
\[
\cos(jx)=\cosh x
\qquad
\sin(jx)=j\sinh x
\]
\[
\cosh(jx)=\cos x
\qquad
\sinh(jx)=j\sin x
\]
\end{schema}

\begin{merke}
Wenn aus einer DGL formal \(\sin(\cdot)\) mit imaginärem Argument entsteht, kannst du sie in \(\sinh/\cosh\) umschreiben.
Das ist in MQS/Wellen-Übergängen praktisch.
\end{merke}

\subsubsection{Inverse Funktionen }
\begin{schema}
\[
\frac{d}{dx}\arctan x = \frac{1}{1+x^2}
\qquad
\int \frac{1}{1+x^2}\,dx = \arctan x + C
\]
\[
\frac{d}{dx}\operatorname{artanh} x = \frac{1}{1-x^2}\quad (|x|<1)
\qquad
\int \frac{1}{1-x^2}\,dx = \operatorname{artanh} x + C
\]
\end{schema}

\subsubsection{Phasor-Form (EM-Standard)}
\begin{schema}
\imp{Realteil-Konvention:}
\[
u(t)=\Re\{\tilde{U}e^{j\omega t}\}
\]
\[
\Re\{Ae^{j\omega t}\}=A\cos(\omega t)
\qquad
\Re\{jAe^{j\omega t}\}=-A\sin(\omega t)
\]
\end{schema}


\begin{checkitem}
\begin{itemize}
  \item Vorzeichenfehler passieren beim Wechsel zwischen \(\sin/\cos\) und komplexen Exponentialformen.
  \item Konsistenzcheck: \(\frac{d}{dt}\cos(\omega t)=-\omega\sin(\omega t)\) muss mit \(j\omega\) im Phasorbereich matchen.
\end{itemize}
\end{checkitem}

\subsubsection{Wann trigonometrisch, wann hyperbolisch?}

\begin{schema}
Ausgangspunkt ist fast immer eine Exponentiallösung:
\[
u(x)=Ae^{\gamma x}+Be^{-\gamma x}
\]
\end{schema}

\begin{schema}
\imp{Umschreibung:}
\[
Ae^{\gamma x}+Be^{-\gamma x}
=
C_1\cosh(\gamma x)+C_2\sinh(\gamma x)
\]
\end{schema}

\begin{checkitem}
\begin{itemize}
  \item \imp{$\gamma \in \mathbb{R}$:} Diffusion, Dämpfung, Skin-Effekt  
        $\Rightarrow$ \(\sinh,\cosh\)
  \item \imp{$\gamma = jk$:} Wellen, Oszillation  
        $\Rightarrow$ \(\sin,\cos\)
\end{itemize}
\end{checkitem}

\begin{merke}
\begin{itemize}
  \item \(\cosh\): gerade Funktion, oft bei Symmetrieachsen.
  \item \(\sinh\): ungerade Funktion, oft bei vorgegebenem Fluss/Feld.
  \item Exponentiell wachsende Terme werden bei offenen Gebieten verworfen.
\end{itemize}
\end{merke}
\subsection{Gradient, Divergenz, Rotation}
% ---------------------------------------------------------

\subsubsection{Gradient}

\[
\nabla \varphi =
\left(
\frac{\partial\varphi}{\partial x},
\frac{\partial\varphi}{\partial y},
\frac{\partial\varphi}{\partial z}
\right)
\]

Anwendung:
\[
\E=-\nabla\varphi
\]

% ---------------------------------------------------------
\subsubsection{Divergenz}

\[
\nabla\cdot\vec F =
\frac{\partial F_x}{\partial x}
+\frac{\partial F_y}{\partial y}
+\frac{\partial F_z}{\partial z}
\]

Anwendung:
\[
\nabla\cdot\J=0
\quad\text{(stationär)}
\]

% ---------------------------------------------------------
\subsubsection{Rotation}

\[
\nabla\times\vec F =
\begin{vmatrix}
\vec e_x & \vec e_y & \vec e_z \\
\partial/\partial x & \partial/\partial y & \partial/\partial z \\
F_x & F_y & F_z
\end{vmatrix}
\]

Anwendung:
\[
\nabla\times\E=-\frac{\partial\B}{\partial t}
\]

% ---------------------------------------------------------
\subsection{Vektoridentitäten }
% ---------------------------------------------------------

\[
\nabla\cdot(\nabla\times\vec A)=0
\]

\[
\nabla\times(\nabla\varphi)=0
\]

\[
\nabla\times(\nabla\times\vec A)
=\nabla(\nabla\cdot\vec A)-\Delta\vec A
\]

\begin{checkitem}
Die letzte Identität ist zentral für die Herleitung der Poisson-Gleichung für $\A$.
\end{checkitem}

% ---------------------------------------------------------
\subsection{Laplace-Operator in 1D-Form}
% ---------------------------------------------------------

Kartesisch:
\[
\Delta u=\frac{\dd^2 u}{\dd x^2}
\]

Zylindrisch (rotationssymmetrisch):
\[
\Delta u=\frac{1}{r}\frac{\dd}{\dd r}\left(r\frac{\dd u}{\dd r}\right)
\]

Sphärisch (kugelsymmetrisch):
\[
\Delta u=\frac{1}{r^2}\frac{\dd}{\dd r}
\left(r^2\frac{\dd u}{\dd r}\right)
\]

% ---------------------------------------------------------
\subsection{Standardlösungen der Laplace-Gleichung}
% ---------------------------------------------------------

\[
\frac{\dd^2 u}{\dd x^2}=0
\Rightarrow u=Ax+B
\]

\[
\frac{1}{r}\frac{\dd}{\dd r}\left(r\frac{\dd u}{\dd r}\right)=0
\Rightarrow u=A\ln r+B
\]

\[
\frac{1}{r^2}\frac{\dd}{\dd r}\left(r^2\frac{\dd u}{\dd r}\right)=0
\Rightarrow u=\frac{A}{r}+B
\]

\begin{merke}
Diese drei Lösungen decken praktisch alle Laplace-Aufgaben in Prüfungen ab.
\end{merke}

% ---------------------------------------------------------
\subsection{Gewöhnliche Differentialgleichungen}
% ---------------------------------------------------------

\subsubsection{Charakteristische Gleichung}

\[
\frac{\dd^2 u}{\dd x^2}+a\frac{\dd u}{\dd x}+bu=0
\]

Ansatz:
\[
u=e^{\lambda x}
\Rightarrow \lambda^2+a\lambda+b=0
\]

% ---------------------------------------------------------
\subsubsection{Spezialfälle}

\[
\lambda=\pm k \Rightarrow u=Ae^{kx}+Be^{-kx}
\]

\[
\lambda=\pm jk \Rightarrow u=A\cos kx+B\sin kx
\]

\begin{schema}
Exponentiell $\Rightarrow$ Diffusion oder Dämpfung.\\
Sinusförmig $\Rightarrow$ Welle oder Resonanz.
\end{schema}

% ---------------------------------------------------------
\subsection{Komplexe Zahlen und Phasoren}
% ---------------------------------------------------------

\[
z=a+jb=|z|e^{j\phi}
\]

\[
|z|=\sqrt{a^2+b^2}
\qquad
\phi=\arctan\left(\frac{b}{a}\right)
\]

\[
e^{j\phi}=\cos\phi+j\sin\phi
\]

% ---------------------------------------------------------
\subsection{Rechnen mit komplexen Größen}
% ---------------------------------------------------------

\[
\frac{1}{a+jb}=\frac{a-jb}{a^2+b^2}
\]

\[
\sqrt{j}=\frac{1+j}{\sqrt{2}}
\]

\begin{checkitem}
Diese Wurzel tritt direkt bei der Skin-Tiefe auf.
\end{checkitem}

% ---------------------------------------------------------
\subsection{Reeller Teil und physikalische Interpretation}
% ---------------------------------------------------------

\[
f(t)=\Re\{\tilde f e^{j\omega t}\}
=|\tilde f|\cos(\omega t+\phi)
\]

\begin{merke}
Komplexe Rechnung ist Rechenhilfe,
physikalisch relevant ist immer der Realteil.
\end{merke}

% ---------------------------------------------------------
\subsection{Integralsätze der Vektorrechnung}

\subsubsection{Gaußscher Satz}

\[
\iiint_V \nabla\cdot\vec F\,\dd V
=
\iint_{\partial V}\vec F\cdot\dd\vec A
\]

\subsubsection{Stokes-Theorem}

\[
\iint_S (\nabla\times\vec F)\cdot\dd\vec A
=
\oint_{\partial S}\vec F\cdot\dd\vec l
\]

\begin{schema}
Diese Sätze verbinden lokale und globale Feldgrößen.
\end{schema}

% ---------------------------------------------------------

\subsection{Delta-Funktion}
% ---------------------------------------------------------

\[
\int_{-\infty}^{\infty}\delta(x-x_0)\dd x=1
\]

\[
\rho(\vec r)=q\,\delta(\vec r-\vec r_0)
\]

Anwendung:
\begin{itemize}
  \item Punktladungen
  \item konzentrierte Quellen
\end{itemize}

% ---------------------------------------------------------
% ---------------------------------------------------------
% 10.X Konzeptionelles Vorgehen: ODE/DGL lösen (homogen + partikulär)
%   Einfügen in Kapitel 10 (Mathematisches Cheatsheet)
%   z.B. nach 10.4 oder vor 10.12
% ---------------------------------------------------------

\subsection{Konzeptionelles Vorgehen: lineare DGL }

\begin{merke}
In Prüfungen ist die Lösungslogik standardisiert:
\[
u = u_{\text{hom}} + u_{\text{part}}
\quad\text{mit}\quad
L[u]=g \Rightarrow
\begin{cases}
L[u_{\text{hom}}]=0\\
L[u_{\text{part}}]=g
\end{cases}
\]
Homogen = Systemstruktur. Partikulär = Antwort auf Quelle. Randbedingungen selektieren die physikalische Lösung.
\end{merke}

\subsubsection{1 Typ erkennen}
\begin{schema}
\begin{itemize}
  \item \imp{Homogen:} $g(x)=0$ \;\;\;(\,keine Quelle\,)
  \item \imp{Inhomogen:} $g(x)\neq 0$ \;\;\;(\,Quelle/Anregung vorhanden\,)
  \item \imp{Linear (Standardform):} $a_n u^{(n)} + \dots + a_1 u' + a_0 u = g(x)$
\end{itemize}
\end{schema}

\subsubsection{2 Homogene Lösung $u_{\text{hom}}$ (charakteristische Gleichung)}
\begin{vorgehen}
\begin{enumerate}
  \item Setze \imp{homogene Gleichung} auf:
  \[
  a_n u^{(n)} + \dots + a_1 u' + a_0 u = 0
  \]
  \item Ansatz (konstante Koeffizienten):
  \[
  u = e^{\lambda x}
  \]
  \item \imp{Charakteristische Gleichung} $P(\lambda)=0$ lösen.
\end{enumerate}
\end{vorgehen}

\begin{schema}
\imp{Typische Fälle (2. Ordnung):}
\begin{itemize}
  \item $\lambda_{1,2}\in\mathbb{R}$:
  \[
  u_{\text{hom}} = C_1 e^{\lambda_1 x} + C_2 e^{\lambda_2 x}
  \]
  \item $\lambda=\alpha \pm j\beta$:
  \[
  u_{\text{hom}} = e^{\alpha x}\left(C_1 \cos(\beta x) + C_2 \sin(\beta x)\right)
  \]
\end{itemize}
\end{schema}

\begin{checkitem}
\begin{itemize}
  \item \imp{Physikalische Selektion:} wachsende Terme im Unendlichen verwerfen (falls gefordert).
  \item Symmetriebedingungen zählen als Randbedingungen.
\end{itemize}
\end{checkitem}

\subsubsection{3 Partikuläre Lösung $u_{\text{part}}$ (nur wenn $g\neq 0$)}
\begin{vorgehen}
\begin{enumerate}
  \item Wähle Ansatzform \imp{entsprechend} Quelle $g(x)$.
  \item Bestimme Koeffizienten durch Einsetzen in die inhomogene DGL.
  \item \imp{Resonanzregel:} Wenn der Ansatz Teil von $u_{\text{hom}}$ ist, mit $x^s$ multiplizieren, bis linear unabhängig.
\end{enumerate}
\end{vorgehen}

\begin{schema}
\imp{Ansatz-Leitfaden (konstante Koeffizienten):}
\begin{itemize}
  \item $g(x)=K$ \;\;\;$\Rightarrow$ Ansatz $u_{\text{part}}=A$
  \item $g(x)=\sum_{k=0}^m b_k x^k$ \;\;\;$\Rightarrow$ Ansatz Polynom Grad $m$
  \item $g(x)=e^{ax}$ \;\;\;$\Rightarrow$ Ansatz $u_{\text{part}}=A e^{ax}$
  \item $g(x)=\cos(\omega x),\sin(\omega x)$ \;\;\;$\Rightarrow$ Ansatz $u_{\text{part}}=A\cos(\omega x)+B\sin(\omega x)$
  \item $g(x)=e^{ax}\cos(\omega x)$ oder $e^{ax}\sin(\omega x)$ \;\;\;$\Rightarrow$ Ansatz $u_{\text{part}}=e^{ax}(A\cos(\omega x)+B\sin(\omega x))$
\end{itemize}
\end{schema}

\subsubsection{4 Gesamtlösung und Konstanten}
\begin{schema}
\[
u(x) = u_{\text{hom}}(x) + u_{\text{part}}(x)
\]
\end{schema}

\begin{vorgehen}
Konstanten $C_i$ bestimmen über:
\begin{itemize}
  \item Randbedingungen (Dirichlet/Neumann/Robin),
  \item Symmetriebedingungen,
  \item Verhalten im Unendlichen (Abklingen/Endlichkeit).
\end{itemize}
\end{vorgehen}

\begin{checkitem}
\begin{itemize}
  \item Dimensionscheck und Grenzfallcheck durchführen.
  \item Stetigkeit an Übergängen prüfen (falls Materialgrenzen).
\end{itemize}
\end{checkitem}

\subsection{Grenzwerte und Plausibilitätschecks}
% ---------------------------------------------------------

\[
\lim_{r\to\infty}\frac{1}{r}=0
\qquad
\lim_{\omega\to 0} j\omega = 0
\]

\begin{merke}
Grenzwerte sind der schnellste Weg,
mathematische Fehler in EM-Aufgaben zu finden.
\end{merke}

% ---------------------------------------------------------

\subsection{Finale Mathematik-Checkliste}
% ---------------------------------------------------------

\begin{checkitem}
\begin{itemize}
  \item Koordinatensystem korrekt?
  \item Ableitungen vollständig?
  \item Logarithmen richtig behandelt?
  \item Komplexe Terme korrekt interpretiert?
  \item Exponentielle Lösungen physikalisch selektiert?
\end{itemize}
\end{checkitem}

\begin{taktik}
Wenn du bei einer EM-Aufgabe festhängst:
Reduziere sie auf Mathematik.
Fast immer ist es eine Laplace-, Poisson-, Diffusions- oder Wellengleichung.
\end{taktik}
% =========================================================

\section{1-seitige Entscheidungs-Matrix für EM-Prüfungen}
% =========================================================
\vspace{4pt}

% ---------------------------------------------------------
\subsection*{Schritt 1: Zeitverhalten}
% ---------------------------------------------------------

\begin{schema}
\begin{tabular}{p{4cm} p{10cm}}
\imp{Frage} & \imp{Konsequenz} \\ \hline
Keine Zeitabhängigkeit & $\partial/\partial t=0$ \\
Sinusförmig, Frequenz $f$ gegeben & Frequenzbereich, $e^{j\omega t}$ \\
Ausbreitung, Phasenverschiebung im Raum & Wellenregime \\
\end{tabular}
\end{schema}

\vspace{2pt}

% ---------------------------------------------------------
\subsection*{Schritt 2: Dominante Physik}
% ---------------------------------------------------------

\begin{schema}
\begin{tabular}{p{4cm} p{10cm}}
\imp{Beobachtung} & \imp{Regime} \\ \hline
Ladungen, Potentiale, Dielektrikum & Elektrostatik \\
Ströme, Widerstand, DC & Stationäre Strömung \\
Ströme treiben $\B$, keine Zeitabhängigkeit & Magnetostatik \\
AC, Leiter, Skin-Effekt & Magneto-Quasistatik (MQS) \\
Reflexion, Transmission, Impedanz & Elektrodynamik (Wellen) \\
\end{tabular}
\end{schema}

\vspace{2pt}

% ---------------------------------------------------------
\subsection*{Schritt 3: Darf ein Potential verwendet werden?}
% ---------------------------------------------------------

\begin{schema}
\begin{tabular}{p{4cm} p{10cm}}
\imp{Bedingung} & \imp{Wahl} \\ \hline
$\nabla\times\E=0$ & $\E=-\nabla\varphi$ \\
$\nabla\cdot\B=0$ & $\B=\nabla\times\A$ \\
Induktion dominant & $\E=-\partial\A/\partial t$ \\
Wellen & direkt $\E,\Hf$ ansetzen \\
\end{tabular}
\end{schema}

\vspace{2pt}

% ---------------------------------------------------------
\subsection*{Schritt 4: Führende Gleichung}
% ---------------------------------------------------------

\begin{schema}
\begin{tabular}{p{5cm} p{9cm}}
\imp{Situation} & \imp{Grundgleichung} \\ \hline
Elektrostatik, $\rho=0$ & $\Delta\varphi=0$ (Laplace) \\
Elektrostatik, $\rho\neq0$ & $\Delta\varphi=-\rho/\eps$ (Poisson) \\
Stationäre Strömung & $\nabla\cdot(\sigma\nabla\varphi)=0$ \\
Magnetostatik & $\Delta\A=-\mu\J$ \\
MQS (Leiter) & $\Delta\A-j\omega\mu\sigma\A=0$ \\
Welle & $\Delta\E+\omega^2\mu\eps\E=0$ \\
\end{tabular}
\end{schema}

\vspace{2pt}

% ---------------------------------------------------------
\subsection*{Schritt 5: Symmetrie $\rightarrow$ Koordinaten}
% ---------------------------------------------------------

\begin{schema}
\begin{tabular}{p{4cm} p{10cm}}
\imp{Geometrie} & \imp{Koordinaten + Abhängigkeit} \\ \hline
Ebene, Platte & kartesisch, $u(x)$ \\
Zylinder, Draht & zylindrisch, $u(r)$ \\
Kugel, Punktquelle & sphärisch, $u(r)$ \\
Unendlich lang & $\partial/\partial z=0$ \\
Rotationssymmetrie & $\partial/\partial\varphi=0$ \\
\end{tabular}
\end{schema}

\vspace{2pt}

% ---------------------------------------------------------
\subsection*{Schritt 6: Standardlösungen merken}
% ---------------------------------------------------------

\begin{schema}
\begin{tabular}{p{5cm} p{9cm}}
\imp{Gleichung} & \imp{Lösungsform} \\ \hline
$\frac{\dd^2 u}{\dd x^2}=0$ & $u=Ax+B$ \\
$\frac{1}{r}\frac{\dd}{\dd r}(r u')=0$ & $u=A\ln r+B$ \\
$\frac{1}{r^2}\frac{\dd}{\dd r}(r^2 u')=0$ & $u=A/r+B$ \\
$u''-\alpha^2 u=0$ & $u=Ae^{\alpha x}+Be^{-\alpha x}$ \\
$u''+k^2u=0$ & $u=A\cos kx+B\sin kx$ \\
\end{tabular}
\end{schema}

\vspace{2pt}

% ---------------------------------------------------------
\subsection*{Schritt 7: Randbedingungen (physikalisch)}
% ---------------------------------------------------------

\begin{schema}
\begin{tabular}{p{5cm} p{9cm}}
\imp{Rand} & \imp{Bedingung} \\ \hline
Idealer Leiter (ES) & $\varphi=\text{konstant}$ \\
Isolator (DC) & $\J_n=0$ \\
Materialgrenze & Stetigkeit tangentialer Felder \\
Symmetrieebene & $\partial u/\partial n=0$ \\
Unendlichkeit & Lösung endlich / abklingend \\
\end{tabular}
\end{schema}

\vspace{2pt}

% ---------------------------------------------------------
\subsection*{Schritt 8: Pflicht-Checks}
% ---------------------------------------------------------

\begin{checkitem}
\begin{itemize}
  \item Dimensionen korrekt?
  \item Grenzfälle: $\omega\to0$, $r\to\infty$, $\sigma\to0$?
  \item Richtungen von $\E$, $\Hf$, $\vec S$ plausibel?
  \item Energieerhaltung bei Wellen erfüllt?
\end{itemize}
\end{checkitem}

\begin{taktik}
Wenn du unsicher bist:  
\imp{Gehe exakt diese 8 Schritte durch.}  
Keine Abkürzungen. Jede EM-Prüfungsaufgabe ist darin vollständig enthalten.
\end{taktik}

\end{document}
